% =========================================================
% דוח מכין – מצפן ספקים
% הקובץ הראשי: מחבר את הפרקים. כל העיצוב נמצא ב-template.cls.
% קומפילציה: lualatex main.tex  (פעמיים-שלוש, לתוכן העניינים)
% =========================================================
\documentclass{template}

\begin{document}

% ---------- עמודים מקדימים ----------
% כל חלק נמצא בתיקייה משלו
% מספור: עמודים פותחים בספרות רומיות, גוף הדוח מתחיל מ-1
% ---------- עמוד שער (ללא מספר עמוד) ----------
% שימו לב: לפי המדריך, מספרי ת"ז מופיעים רק בעותק שמוגש למנחה ולוועדה.
% אם המאגר ב-GitHub ציבורי – לא להכניס כאן ת"ז.
\begin{titlepage}
\centering
\setlength{\parskip}{0pt}
\vspace*{-1.5cm}

% לוגואים: הקבצים נמצאים ב"עמוד שער ותקציר/logos/".
% אם קובץ חסר, מוצגת מסגרת עם טקסט במקומו.
\newcommand{\logo}[3]{\IfFileExists{#1}{\includegraphics[height=#2]{#1}}{\fbox{\parbox[c][1.6cm][c]{4cm}{\centering #3}}}}
\logo{../עמוד שער ותקציר/logos/biu-logo.png}{1.8cm}{[לוגו בר-אילן]}
\hfill
\logo{../עמוד שער ותקציר/logos/decixai-logo.png}{1.5cm}{[לוגו הצוות]}
\hfill
\logo{../עמוד שער ותקציר/logos/iai-logo.png}{1.6cm}{[לוגו התעשייה האווירית]}

\vspace{1cm}
{\large אוניברסיטת בר-אילן -- המחלקה לניהול, מסלול ניהול טכנולוגיות\par}
\vspace{0.2cm}
{\large סמינריון מחקרי בניהול טכנולוגיות (5523001)\par}

\vspace{1.2cm}
{\LARGE\bfseries הערכת ספקים בשלב טרום ההתקשרות מתוך מידע ציבורי\\[0.2cm]
באמצעות מערכת סוכני בינה מלאכותית:\\[0.2cm]
מתודולוגיה לחברות ציבוריות ופרטיות בתעשייה הביטחונית-אווירית\par}

\vspace{0.8cm}
{\large\bfseries\foreignlanguage{english}{Pre-Contract Supplier Assessment from Public Information\\
Using a Multi-Agent AI System:\\
A Methodology for Public and Private Companies\\
in the Defense-Aerospace Industry}\par}

\vspace{1cm}
{\Large דוח מכין\par}

\vspace{1cm}
% בעברית (bidi=basic) העמודות בטבלה נקראות משמאל לימין, לכן התווית בעמודה השנייה
\begin{tabular}{r >{\bfseries}r}
טל אזואלוס & מגישות: \\
עמית סוסן & \\[0.3cm]
ד"ר יורם סגל, אוניברסיטת בר-אילן & מנחה אקדמי: \\[0.3cm]
מר מירום לוין, התעשייה האווירית לישראל & מנחה תעשייתי: \\[0.3cm]
25.11.2026 & תאריך הגשה: \\
\end{tabular}

\vfill
\end{titlepage}

\pagenumbering{roman}
\setcounter{page}{1}
% ---------- תקציר מנהלים ----------
% 22.11.2026: גוף שלישי, עד 200 מילים (הנחיית המרצה). כל מספר מגובה במקור בפרק 1.
\chapter*{תקציר}
\addcontentsline{toc}{chapter}{תקציר}

חברה ביטחונית-אווירית צריכה ספקים שמייצרים מוצרים איכותיים ואמינים, כי ספק חלש עלול לפגוע ברציפות האספקה של מערכות להגנת המדינה. בארגונים גדולים, קליטת ספק חדש נמשכת שבועות. הקושי גדול במיוחד בחברות פרטיות: שיעור הסיכון הפיננסי בהן גבוה יותר, ודווקא עליהן כמעט אין דוחות כספיים פומביים.

הדוח מציע מתודולוגיה להערכת ספק לפני ההתקשרות, שנשענת על מידע ציבורי בלבד ומבוצעת בידי סוכני בינה מלאכותית. ההערכה בנויה משני חלקים: \textbf{פרופיל חוזקות} (איכות ואמינות, התאמה טכנית, עמידה ביעדים ומצב פיננסי) ו\textbf{דגלי סיכון} (מצב גיאופוליטי, אירועים חריגים ואבטחת סייבר). ספק ברשימת סנקציות או ממדינה עוינת לישראל נפסל לפני כל בדיקה אחרת. המערכת מציגה גם כמה מהמידע הנדרש נמצא, וכשהוא חסר היא מפיקה רשימה לבדיקה ידנית. התוצר הוא המלצה מנומקת עם מקור לכל טענה, וההחלטה הסופית נשארת בידי אדם.

אבני הדרך: אב-טיפוס עובד בשבוע 10 של תוכנית העבודה, ותיקוף לאחור על מקרי כשל ציבוריים ועל ספקים מוכחים עד שבוע 13.

\vspace{0.5cm}
\noindent\textbf{מילות מפתח:} ניהול סיכוני ספקים; הערכת ספקים טרום התקשרות; מערכות מרובות סוכנים; מידע ממקורות גלויים (\foreignlanguage{english}{OSINT}); אבטחת סייבר בשרשרת האספקה; תעשייה ביטחונית-אווירית; מידע חסר.

\clearpage

% תוכן העניינים והרשימות בריווח צפוף יותר
{\setstretch{1.15}
\tableofcontents
\clearpage
\listoffigures
\addcontentsline{toc}{chapter}{רשימת איורים}
\listoftables
\addcontentsline{toc}{chapter}{רשימת טבלאות}
}
\clearpage
% ---------- רשימת קיצורים ומונחים ----------
% כל טקסט באנגלית עטוף ב-\en{...} כדי שלא יתהפך בתוך טבלה בעברית
\chapter*{רשימת קיצורים ומונחים}
\addcontentsline{toc}{chapter}{רשימת קיצורים ומונחים}

{\small
\begin{longtable}{|>{\raggedright\arraybackslash}p{2.2cm}|>{\raggedright\arraybackslash}p{6.2cm}|>{\raggedright\arraybackslash}p{5.2cm}|}
\hline
\textbf{קיצור} & \textbf{פירוש מלא} & \textbf{הסבר} \\ \hline
\endhead
\en{AHP} & \en{Analytic Hierarchy Process} & שיטה לקביעת משקלות בהשוואות בזוגות \\ \hline
\en{AI} & \en{Artificial Intelligence} & בינה מלאכותית \\ \hline
\en{AIDM} & \en{AI-driven Decision-Making} & המסגרת של מאמר הדגל \\ \hline
\en{API} & \en{Application Programming Interface} & ממשק תכנות לגישה אוטומטית למקור מידע \\ \hline
\en{AS9100} & \en{Aerospace Quality Management Standard} & תקן ניהול איכות לתעשיית התעופה \\ \hline
\en{BIS} & \en{Bureau of Industry and Security} & הלשכה האמריקאית לבקרת יצוא \\ \hline
\en{CAGR} & \en{Compound Annual Growth Rate} & שיעור צמיחה שנתי ממוצע \\ \hline
\en{CMMC} & \en{Cybersecurity Maturity Model Certification} & הסמכת סייבר לספקי משרד ההגנה האמריקאי \\ \hline
\en{CR} & \en{Consistency Ratio} & יחס העקביות של ההשוואות ב-\en{AHP} \\ \hline
\en{ENISA} & \en{European Union Agency for Cybersecurity} & סוכנות הסייבר של האיחוד האירופי \\ \hline
\en{ESG} & \en{Environmental, Social, Governance} & מדדי סביבה, חברה וממשל תאגידי \\ \hline
\en{FHR} & \en{Financial Health Rating} & ציון בריאות פיננסית של \en{RapidRatings} \\ \hline
\en{FIR} & \en{False Identification Rate} & שיעור זיהוי שגוי (המאמר הנלווה) \\ \hline
\en{GAO} & \en{Government Accountability Office} & משרד מבקר המדינה האמריקאי \\ \hline
\en{GDELT} & \en{Global Database of Events, Language, and Tone} & מאגר חדשות עולמי פתוח \\ \hline
\en{GPR} & \en{Geopolitical Risk Index} & מדד סיכון גיאופוליטי מבוסס עיתונות \\ \hline
\en{IAQG} & \en{International Aerospace Quality Group} & הגוף שמנהל את הסמכות \en{AS9100} \\ \hline
\en{ISO 9001} & \en{Quality Management Systems Standard} & תקן בינלאומי למערכות ניהול איכות \\ \hline
\en{KYC} & \en{Know Your Customer} & בדיקת "הכר את הלקוח" \\ \hline
\en{LARD-SC} & שם המסגרת במאמר הנלווה \cite{zhao2025} & זיהוי סיכוני ספקים מחדשות בעזרת מודל שפה \\ \hline
\en{LLM} & \en{Large Language Model} & מודל שפה גדול \\ \hline
\en{MCDM} & \en{Multi-Criteria Decision Making} & קבלת החלטות רב-קריטריונית \\ \hline
\en{NIST} & \en{National Institute of Standards and Technology} & מכון התקנים האמריקאי \\ \hline
\en{OFAC} & \en{Office of Foreign Assets Control} & משרד הסנקציות של משרד האוצר האמריקאי \\ \hline
\en{OSINT} & \en{Open-Source Intelligence} & מידע ממקורות גלויים \\ \hline
\en{PRR} & \en{Potential Risk Rate} & שיעור הסיכונים שאושרו חלקית (המאמר הנלווה) \\ \hline
\en{QFD} & \en{Quality Function Deployment} & פריסת פונקציית האיכות \\ \hline
\en{RACI} & \en{Responsible, Accountable, Consulted, Informed} & מטריצת אחריות \\ \hline
\en{RAG} & \en{Retrieval-Augmented Generation} & יצירת טקסט שמבוססת על אחזור מקורות \\ \hline
\en{RPI} & \en{Relative Performance Index} & מדד ביצוע יחסי (המאמר הנלווה) \\ \hline
\en{RVR} & \en{Risk Validation Rate} & שיעור אימות מלא (המאמר הנלווה) \\ \hline
\en{SAM} & \en{Serviceable Available Market} & הפלח שאליו פונה הפתרון \\ \hline
\en{SDN} & \en{Specially Designated Nationals} & רשימת הסנקציות של משרד האוצר האמריקאי \\ \hline
\en{SEC} & \en{Securities and Exchange Commission} & רשות ניירות הערך האמריקאית \\ \hline
\en{SIPRI} & \en{Stockholm International Peace Research Institute} & מכון המחקר לשלום בשטוקהולם \\ \hline
\en{SMART} & \en{Specific, Measurable, Achievable, Relevant, Time-bound} & עקרונות לניסוח יעדים \\ \hline
\en{SOM} & \en{Serviceable Obtainable Market} & הנתח שאפשר לתפוס בפועל \\ \hline
\en{TAM} & \en{Total Addressable Market} & כל השוק העולמי בתחום \\ \hline
\en{WBS} & \en{Work Breakdown Structure} & מבנה פירוק העבודה \\ \hline
\en{WGI} & \en{Worldwide Governance Indicators} & מדדי הממשל של הבנק העולמי \\ \hline
\en{WP} & \en{Work Package} & חבילת עבודה \\* \hline
\end{longtable}
}

\clearpage


% ---------- גוף הדוח (מספור רגיל מ-1) ----------
\clearpage
\pagenumbering{arabic}
\setcounter{page}{1}
% ---------- פרק 1: מבוא ורקע ----------
% 22.11.2026: גוף שלישי וקיצור. כל מספר מגובה במקור מ"מקורות/references.tex".
% מבנה משפך: מגמה עולמית ← הארגון ← המוטיבציה ← מבנה הדוח.
\chapter{מבוא ורקע}

\section{המגמה: סיכוני שרשרת האספקה וכניסת הבינה המלאכותית}

שרשראות האספקה של תעשיית התעופה והביטחון נמצאות תחת לחץ גובר: הביקוש למערכות ביטחוניות עולה, והחוסן הפיננסי של רבים מהספקים נשחק. מבחן לחץ של חברת הדירוג \en{RapidRatings} מ-2026 בדק 4,309 ספקים בענף. במצב הנוכחי 17.6\% מהם נמצאים בסיכון פיננסי גבוה או גבוה מאוד, ובתרחיש הלחץ השיעור עולה ל-20.3\%. בחברות הפרטיות השיעור הנוכחי הוא 22.1\%, לעומת 15.9\% בציבוריות, ובתרחיש הלחץ 26.7\% ו-17.9\% בהתאמה \cite{rapidratings2026}. הנתונים מגיעים מחברה שמוכרת דירוגי ספקים ולא ממחקר שפיט, ולא נמצא להם תחליף אקדמי.

הקושי אינו רק בסיכון עצמו, אלא ביכולת לראות אותו. משרד מבקר המדינה האמריקאי (\en{GAO}) מצא ב-2025 שבשרשרת של מטוס ה-\en{F-35} יש בידי משרד ההגנה מידע על פחות מ-10\% מספקי הרכיבים וחומרי הגלם, ובבדיקה נוספת 22 מתוך 63 ספקים שהתבקשו למסור נתונים על מדינת המקור סירבו או לא השיבו \cite{gao2025}. כלומר, גם ארגון בעל משאבים רבים אינו יכול להסתמך רק על מידע שהספק מוסר מרצונו.

ממד נוסף הוא סיכון הסייבר דרך הספקים. סוכנות הסייבר האירופית (\en{ENISA}) ניתחה 24 מתקפות על שרשראות אספקה, ומצאה שבכ-66\% מהמקרים הספקים לא ידעו או לא דיווחו כיצד נפרצו \cite{enisa2021}. ב-2019 עצרה מתקפת כופרה את הייצור במפעלי \en{ASCO Industries}, יצרנית בלגית של חלקי מבנה למטוסים, בכמה מדינות \cite{certeu2019}. גם הרגולציה מגיבה: משרד ההגנה האמריקאי קבע תקן מחייב לאבטחת סייבר אצל ספקיו (\en{CMMC}), שנכנס לתוקף בסוף 2024 \cite{cmmc2024}.

במקביל מתפתחים כלים מבוססי בינה מלאכותית: מודלי שפה גדולים (\en{LLM}) שמזהים סיכונים מתוך כתבות על ספקים \cite{zhao2025}, ומערכות של כמה סוכני \en{AI} שמקשרות אירועים לרשת הספקים של הארגון \cite{almahri2026}. רוב העבודות האלה עוסקות בניטור ספקים קיימים, ולא בהערכת ספק חדש לפני ההתקשרות (פרק~\ref{ch:literature}).

\section{ההקשר היישומי: התעשייה האווירית ותהליך טרום ההתקשרות}

התעשייה האווירית לישראל (תע"א) היא חברה ממשלתית לפי חוק החברות הממשלתיות, ובסוף 2022 עבדה עם כ-5,100 ספקים \cite{mevaker2024}. היא מפתחת ומייצרת מערכות אוויריות, חלליות וביטחוניות. ב-2025 הגיעו הכנסותיה לכ-7.38 מיליארד דולר, כ-66\% מהן מייצוא, וצבר ההזמנות עמד על כ-29 מיליארד דולר \cite{iai2026}.

לפני כל התקשרות עם ספק חדש צריך לענות על שתי שאלות. הראשונה היא \textbf{האם הספק טוב}: האם הוא מייצר מוצרים איכותיים ואמינים, מסוגל לייצר לפי מפרט הרכיב, עומד ביעדים שהציב ויציב כלכלית. השנייה היא \textbf{האם יש סיבה לחשוש ממנו}: סנקציות ומגבלות יצוא, אבטחת הסייבר שלו, יציבות המדינה שבה הוא פועל ויחסיה עם ישראל. מאחר שתע"א עוסקת במערכות להגנת המדינה ובמידע מסווג, ספק לא מתאים עלול לפגוע ברציפות האספקה, באיכות המערכות ובביטחון המידע.

\section{המוטיבציה: למה התהליך הקיים כבר לא מספיק}

בסקר של חברת \en{HICX} מ-2021, בקרב 100 מנהלי רכש בכירים בחברות שמחזורן מיליארד דולר ומעלה, 63\% דיווחו שקליטת ספק חדש נמשכת 15 עד 28 ימים, ו-31\% דיווחו על 29 עד 60 ימים \cite{hicx2021}. הסקר עסק בקליטת ספקים באופן כללי, ולא נמצא נתון פומבי על משך הבדיקה בתעשייה ביטחונית.

הבעיה חריפה במיוחד בחברות פרטיות: שיעור הסיכון הפיננסי בהן גבוה יותר \cite{rapidratings2026}, ודווקא עליהן כמעט אין דוחות כספיים פומביים. \en{Incora}, מפיצה פרטית של חלקים לתעשיית התעופה שסיפקה ליותר מ-8,400 לקוחות, פשטה רגל ב-2023, אחרי ירידה של 18.3\% בהכנסות ונטל חוב כבד \cite{incora2023}.

מכאן הצורך בגישה שתהיה מהירה ושיטתית, תישען על מידע ציבורי שאפשר לאמת, תעריך את \textbf{חוזקות} הספק ולא רק תחפש סיכונים, תזהה \textbf{דגלי סיכון} מכל סוג ותדרג אותם לפי ההשפעה על רציפות האספקה, ותציג במפורש מה לא נמצא ומה צריך לבדוק ידנית. הכלי המוצע בדוח, \textbf{מצפן ספקים}, מפיק הערכה מלאה ומנומקת של הספק, שמובילה להמלצה אם להתקשר איתו.

\section{מבנה הדוח}

פרק~\ref{ch:problem} מגדיר את הבעיה, את שאלת המחקר, את ההשערות ואת מפרט הדרישות. פרק~\ref{ch:literature} סוקר את הספרות ומגדיר את פער הידע. פרק~\ref{ch:market} מציג את השוק ואת המוצרים הקיימים. פרק~\ref{ch:design} בוחר בין חלופות תכן ומתאר את הארכיטקטורה, את מקורות המידע ואת החדשנות. פרק~\ref{ch:workplan} מציג את תוכנית העבודה, פרק~\ref{ch:risks} את הסיכונים, ופרק~\ref{ch:summary} מסכם את המוכנות לשלב המימוש.

% ---------- פרק 2: הגדרת הבעיה, שאלת המחקר והיעדים ----------
% 22.11.2026: גוף שלישי וקיצור. ערכי היעד (90%, שעה, 4 מתוך 5) אושרו ע"י עמית 13.11; 80% ב-REQ-PER-03 והגדרת ספק מוכח – הצעות Claude.
\chapter{הגדרת הבעיה, שאלת המחקר והיעדים}\label{ch:problem}

\section{המצב הרצוי}

הבעיה היא \textbf{איכות ההחלטה ומשך הבדיקה} בשלב טרום ההתקשרות עם ספק חדש בתעשייה הביטחונית-אווירית. טבלה~\ref{tab:goal} מגדירה את המצב הרצוי בחמישה ממדים.

{\small
\begin{longtable}{|>{\raggedright\arraybackslash}p{3.2cm}|>{\raggedright\arraybackslash}p{10.8cm}|}
\caption{המצב הרצוי בסוף הפרויקט}\label{tab:goal}\\
\hline
\textbf{ממד} & \textbf{המצב הרצוי} \\ \hline
\endfirsthead
\hline
\textbf{ממד} & \textbf{המצב הרצוי} \\ \hline
\endhead
משך ההערכה & הערכה מלאה ומנומקת מתוך מידע ציבורי בפחות משעה; מה שאפשר לאמת רק אצל הספק מסומן לבדיקה ידנית \\ \hline
איכות ואמינות & פרופיל חוזקות מאותות ציבוריים: הסמכות, לקוחות וחוזים מוכרים, הכרה פומבית והיסטוריית ריקולים \\ \hline
חברות פרטיות & הערכה גם בלי דוחות כספיים, בעזרת מדדים עקיפים, עם ציון גלוי של כמה מהמידע נמצא \\ \hline
שקיפות & מקור לכל טענה, ורשימה מפורשת של מה שלא נמצא \\ \hline
דגלי סיכון & זיהוי שיטתי של אירועים מכל סוג, ודירוגם לפי ההשפעה על רציפות האספקה \\ \hline
\end{longtable}
}

\textbf{הפער ומחירו.} בארגונים גדולים, קליטת ספק חדש נמשכת שבועות \cite{hicx2021}, וגם בסופה אין מדד שקוף לשאלה כמה מההחלטה נשען על מידע חסר. ספק חדש מוערך לפני שיש איתו היסטוריית עבודה, ולכן קשה לדעת אם מוצריו איכותיים. הפער חמור במיוחד בחברות פרטיות. \en{Incora}, מפיצה פרטית של חלקי תעופה, פשטה רגל ב-2023 אחרי ירידה בהכנסות, נטל חוב כבד ועמידה בזמני האספקה בכ-50\% בלבד מהמשלוחים \cite{incora2023}; \en{ASCO Industries} הותקפה ב-2019 בכופרה, והייצור שלה נעצר \cite{certeu2019}. בשני המקרים נפגעה רציפות האספקה ללקוחות.

\section{שאלת המחקר וההשערות}

הפרויקט משלב \textbf{שאלת מחקר}, שנבדקת דרך השערות שאפשר להפריך, ו\textbf{יעדי תכן} לפי עקרונות \en{SMART}, שמגדירים את אב-הטיפוס.

\begin{quote}
\textbf{שאלת המחקר: באיזו מידה} מתודולוגיה מבוססת סוכני בינה מלאכותית, שנשענת \textbf{על מידע ציבורי בלבד}, מצליחה \textbf{להעריך את האיכות, האמינות והסיכונים} של ספק ולהפיק המלצת התקשרות מנומקת ושקופה \textbf{בשלב טרום ההתקשרות} בתעשייה הביטחונית-אווירית?
\end{quote}

\textbf{שאלות משנה:}
\begin{enumerate}
\item אילו מדדים ציבוריים בכל אחד משבעת הממדים קשורים בספרות לאיכות ספק או לכשל שלו? ארבעה ממדים בונים את \textbf{פרופיל החוזקות} (איכות ואמינות, התאמה טכנית, עמידה ביעדים ומצב פיננסי), ושלושה את \textbf{דגלי הסיכון} (מצב גיאופוליטי, אירועים חריגים ואבטחת סייבר).
\item על חברה פרטית חסרים בדרך כלל דוחות כספיים, יעדים שפורסמו, דיווחי סייבר (חובה שחלה על חברות ציבוריות בארה"ב \cite{sec2023}) ושקיפות בעלות. אילו מדדים עקיפים זמינים לה, ועד כמה הם מחליפים את המידע החסר?
\item איך לשקלל את הממדים כשחלק מהמידע חסר, ואיך החוסר משפיע על ההמלצה ועל רמת הביטחון בה?
\item האם המתודולוגיה מזהה מראש מקרי כשל ציבוריים, וכמה זמן לפני מועד הכשל?
\end{enumerate}

\textbf{הגדרה אופרטיבית של כשל ספק.} כשל נמדד לפי ההשפעה על רציפות האספקה, ולא לפי גודל ההפסד, בסולם של ארבע רמות:
\begin{enumerate}
\item \textbf{חוסם:} הספק ברשימת סנקציות או בקרת יצוא, או שהוא או בעלי השליטה בו פועלים ממדינה עוינת לישראל. בכלי נקבע שמדינה עוינת היא מדינה שישראל מגדירה כמדינת אויב, או מדינה שאין לה יחסים דיפלומטיים עם ישראל. רשימת המדינות תיקבע מול מקור רשמי בשלב בניית השער (\en{T2.1}).
\item \textbf{כשל אספקה:} אירוע שעצר בפועל ייצור או אספקה, כמו פשיטת רגל, אובדן הסמכה, השבתה בגלל מתקפת סייבר, או מלחמה במדינת הספק.
\item \textbf{אות סיכון:} אירוע שמחליש את הספק בלי עדות לעצירת אספקה, כמו הונאה, תביעה, החלפת הנהלה, פריצת סייבר שלא עצרה ייצור או הסלמה מדינית.
\item \textbf{רעש:} אזכור שלילי שאינו מהותי.
\end{enumerate}
ספק שהגיע לרמה 1 או 2 הוא \textbf{מקרה כשל} לצורך התיקוף. בגלל המידע המסווג שתע"א עוסקת בו, אירוע סייבר חמור שלא נסגר מונע "מומלץ" עד לבדיקת אבטחה ידנית, גם אם לא עצר ייצור.

\textbf{הגדרה אופרטיבית של ספק מוכח.} ספק שפועל בענף התעופה עשר שנים לפחות בלי אירוע ברמה 1 או 2. ההגדרה נשענת על ותק בלבד, ולא על הסמכות או פרסים, כי אלה אותות שהכלי עצמו מודד, ובחירה לפיהם הייתה הופכת את התיקוף למעגלי.

\textbf{סיווג נכון.} להערכה חמש תוצאות אפשריות (פרק~\ref{ch:design}): מומלץ, מומלץ בתנאים, לא מומלץ, פסילה בשער ומידע לא מספיק. סיווג נכון של ספק שכשל הוא "לא מומלץ" או "פסילה", ושל ספק מוכח "מומלץ" בלבד; "מומלץ בתנאים" נספר כשגיאה בשתי הקבוצות, כדי שכלי שעונה "בתנאים" על כל ספק לא ייראה מדויק. "מידע לא מספיק" מדווח בנפרד, כשיעור הימנעות. בנוסף מדווח ההפרש בציון $m$ בין הקבוצות. שיעור הסיווג הנכון מחושב על מקרי כשל שאפשר לצפות ועל ספקים מוכחים; שאר המקרים במאגר בודקים רכיב מסוים של הכלי (טבלה~\ref{tab:cases}).

\textbf{ההשערות} מוצגות בטבלה~\ref{tab:hypotheses}. מדגם התיקוף קטן, כעשרה עד עשרים מקרים, ולכן הן נבדקות כמקרי בוחן ומדווחות תיאורית, בלי מבחן מובהקות. קו הבסיס הפיננסי של \en{H1} הוא מדד \en{Z} של \en{Altman} \cite{altman1968}. אפשר לחשב אותו רק לחברה שמפרסמת דוחות כספיים, ולכן לחברה פרטית קו הבסיס אינו מחזיר תשובה, וההשוואה מדווחת בנפרד לחברות ציבוריות ולחברות פרטיות.

{\small
\begin{longtable}{|p{0.7cm}|>{\raggedright\arraybackslash}p{4.4cm}|>{\raggedright\arraybackslash}p{4.6cm}|>{\raggedright\arraybackslash}p{4.1cm}|}
\caption{ההשערות, אופן הבדיקה ומה יפריך כל אחת}\label{tab:hypotheses}\\
\hline
\textbf{מזהה} & \textbf{השערה} & \textbf{איך תיבדק} & \textbf{מה יפריך אותה} \\ \hline
\endfirsthead
\hline
\textbf{מזהה} & \textbf{השערה} & \textbf{איך תיבדק} & \textbf{מה יפריך אותה} \\ \hline
\endhead
\en{H1} & הערכה בשני חלקים מבחינה בין ספקים שכשלו לספקים מוכחים טוב יותר מבדיקה פיננסית בלבד & בדיקה לאחור על שתי הקבוצות מול קו בסיס פיננסי; לכל כשל נמדד גם זמן ההתרעה & שיעור הסיווג הנכון אינו גבוה מקו הבסיס, או שהפרש $m$ אינו חיובי \\ \hline
\en{H2} & בחברות פרטיות הכיסוי נמוך יותר, אבל המדדים העקיפים עדיין מבחינים בין הקבוצות & השוואה בין כשלים פרטיים לספקים מוכחים פרטיים & המדדים העקיפים אינם מבחינים בין הקבוצות \\ \hline
\en{H3} & ההסבר, שני החלקים ורשימת "לבדוק ידנית" שימושיים להחלטה & שאלון בסולם 1--5 על דוחות לדוגמה, למנחה התעשייתי ולאנשי רכש נוספים & ציון ממוצע נמוך מ-4 \\ \hline
\end{longtable}
}

\textbf{יעדי תכן (\en{SMART}).} \textbf{ספציפי:} אב-טיפוס שמפיק דוח הערכה לפי שבעת הממדים, עם שער חסימה, שני חלקים, ציון כיסוי והמלצה מדורגת. \textbf{מדיד:} דוח מלא לכל ספק במדגם התיקוף, עם מקור לכל טענה. \textbf{בר-השגה:} מידע ציבורי ומודלי שפה זמינים. \textbf{רלוונטי:} עונה על הפער שלעיל. \textbf{תחום בזמן:} אב-טיפוס בסוף שבוע 10 ותיקוף בסוף שבוע 13 של תוכנית העבודה (פרק~\ref{ch:workplan}).

\section{מפרט הדרישות}

לכל דרישה מוגדרים מזהה, ערך יעד ושיטת אימות (טבלה~\ref{tab:requirements}).

{\small
\begin{longtable}{|>{\raggedright\arraybackslash}p{2.35cm}|>{\raggedright\arraybackslash}p{4.6cm}|>{\raggedright\arraybackslash}p{1.3cm}|>{\raggedright\arraybackslash}p{2.7cm}|>{\raggedright\arraybackslash}p{2.4cm}|}
\caption{מפרט הדרישות}\label{tab:requirements}\\
\hline
\textbf{מזהה} & \textbf{הדרישה} & \textbf{סיווג} & \textbf{ערך יעד} & \textbf{אימות} \\ \hline
\endfirsthead
\hline
\textbf{מזהה} & \textbf{הדרישה} & \textbf{סיווג} & \textbf{ערך יעד} & \textbf{אימות} \\ \hline
\endhead
\en{REQ-FUN-01} & שער חסימה לפני הציון: סנקציות, בקרת יצוא, מדינת מושב ובעלות & תפקודית & 100\% מהישויות שזוהו; בעלות לא ידועה: "לא נבדק" ובדיקה ידנית & ישויות שידוע שהן ברשימות \\ \hline
\en{REQ-FUN-02} & ציון 0--100 ומשקל לכל ממד; משקל של ממד חסר מחולק מחדש באופן גלוי & תפקודית & 7 מתוך 7 ממדים, עם ציון או "אין מספיק מידע" & בדיקת פלט \\ \hline
\en{REQ-FUN-03} & ציון כיסוי מידע לכל ממד ובסך הכול & תפקודית & בכל דוח & בדיקת פלט \\ \hline
\en{REQ-FUN-04} & סיווג זמינות המידע: ציבורית (\en{A}), פרטית חלקית (\en{B}), כמעט ללא מידע (\en{C}) & תפקודית & 90\% התאמה לסיווג ידני & מדגם התיקוף \\ \hline
\en{REQ-FUN-05} & רשימת נושאים לבדיקה ידנית & תפקודית & בכל דוח עם כיסוי חלקי & בדיקת פלט \\ \hline
\en{REQ-FUN-06} & זיהוי פתוח של אירועים שליליים, גם מסוג שאינו ברשימה & תפקודית & לכל אירוע: קטגוריה, חומרה 1--4, ביטחון ומקור & מקרי התיקוף \\ \hline
\en{REQ-FUN-07} & חומרה לפי השפעה על רציפות האספקה & תפקודית & רמה 1: פסילה; רמה 2 שלא נסגרה: "לא מומלץ" & מקרי התיקוף \\ \hline
\en{REQ-FUN-08} & ממד סייבר: הסמכות, אירועים פומביים, ודירוג חיצוני רק אם פורסם & תפקודית & ציון, כיסוי ורשימה ידנית & מקרי סייבר ציבוריים \\ \hline
\en{REQ-FUN-09} & המלצה בשלוש רמות ושתי תוצאות עצירה (פסילה, מידע לא מספיק), עם הסבר לכל ממד ולכל חלק & תפקודית & 100\% מהדוחות & בדיקת פלט \\ \hline
\en{REQ-FUN-10} & ממד איכות: הסמכות (למשל \en{AS9100}), לקוחות וחוזים, ריקולים והוראות תקינות & תפקודית & ציון, כיסוי ורשימה ידנית & הספקים המוכחים \\ \hline
\en{REQ-PER-01} & הבחנה בין ספקים שכשלו לספקים מוכחים & ביצועים & טובה מקו בסיס פיננסי (\en{H1}) & בדיקה לאחור \\ \hline
\en{REQ-PER-02} & זמן הפקת דוח לספק & ביצועים & פחות משעה & מדידה על מדגם התיקוף \\ \hline
\en{REQ-PER-03} & זיהוי האירועים הידועים במקרי התיקוף & ביצועים & 80\% לפחות & מול רשימה ידנית שנבנית מראש \\ \hline
\en{REQ-NFR-01} & מקור לכל טענה עובדתית & הסבריות & 100\% מהטענות & ביקורת על מדגם \\ \hline
\en{REQ-NFR-02} & ההחלטה הסופית בידי אדם & אילוץ & אין החלטה אוטומטית מלבד השער & סקירת תהליך \\ \hline
\en{REQ-CON-01} & מידע ציבורי בלבד & אילוץ & אפס מקורות פנימיים & ביקורת מקורות \\ \hline
\en{REQ-CON-02} & איסוף פסיבי, בלי סריקת מערכות הספק & אילוץ חוקי & אפס סריקות & סקירת קוד \\ \hline
\en{REQ-CON-03} & אפשרות להריץ בעתיד על מודל מקומי & אילוץ אבטחה & בלי תלות בספק ענן & סקירת ארכיטקטורה \\ \hline
\end{longtable}
}

% ---------- פרק 3: סקירת הספרות ומאמר הדגל ----------
% 22.11.2026: גוף שלישי וקיצור. מבנה תמטי וביקורתי, לפי פרק 5 במדריך.
% כל מקור נבדק מול DOI/arXiv או טקסט מלא. סטטוס: "מקורות - רשימה מאומתת.md".
% מאמר הדגל (Sadeghi 2026) והמאמר הנלווה (Zhao 2025) נקראו בטקסט מלא; כל המספרים בסעיפים שלהם לקוחים מהם.
\chapter{סקירת הספרות ומאמר הדגל}\label{ch:literature}

הסקירה מחולקת לפי תחומי ידע, וכל תחום מתאים לחלק אחר בבעיה (פרק~\ref{ch:problem}): הערכה רב-קריטריונית וקביעת משקלות, איכות ספקים, סיכון גיאופוליטי, כשל פיננסי, זיהוי אירועים שליליים, אבטחת סייבר ומודלי שפה בשרשרת אספקה. בסוף הפרק מוצגות מטריצת הספרות, מאמר הדגל, המאמר הנלווה ופער הידע.

% =====================================================================
\section{הערכת ספקים וקבלת החלטות רב-קריטריונית}\label{sec:lit-mcdm}

\en{Ho}, \en{Xu} ו-\en{Dey} סקרו את שיטות קבלת ההחלטות הרב-קריטריוניות (\en{MCDM}) לבחירת ספקים, והסיקו שהן עדיפות על בחירה לפי עלות בלבד \cite{ho2010}. לכן ההמלצה המוצעת נשענת על כמה ממדים בעלי משקלות. השיטה הנפוצה לקביעת משקלות היא תהליך ההיררכיה האנליטית (\en{AHP}) של \en{Saaty} \cite{saaty1990}: ההחלטה מפורקת להיררכיה של מטרה, קריטריונים ותתי-קריטריונים, כל זוג קריטריונים מושווה, והמשקלות מחושבים מההשוואות, תוך בדיקת \textbf{עקביות} השיפוטים. המבנה ההיררכי מתאים למודל המשקלות המוצע (חלק, ממד וקטגוריות, פרק~\ref{ch:design}). \en{Boukrouh} ועמיתיו השוו שלוש דרכים לקביעת משקלות לבחירת ספקים: שיפוט (\en{AHP}), חישוב מהנתונים (אנטרופיה) ושילוב של השתיים \cite{boukrouh2024}. בניתוח רגישות על 21 תרחישים, השילוב נשאר יציב ב-19, החישוב ב-3, והשיפוט בלבד באף אחד.

\textbf{ביקורת.} לשיטות הקלאסיות שתי הנחות שלא מתקיימות כאן. הראשונה היא \textbf{מידע מלא}: השיטות מניחות ציון לכל ספק בכל קריטריון, ואינן אומרות מה לעשות כשלחברה פרטית אין דוחות כספיים או כיסוי בחדשות; בזה עוסקת שאלת המשנה השלישית. השנייה נוגעת ל\textbf{קביעת המשקלות}: בפרויקט זה את ההשוואות בזוגות ממלאות מגישות הדוח, כל השוואה מנומקת במקור מהספרות, והמשקלות מוצגים למנחה התעשייתי להערות לפני שהם נקבעים (\en{T1.3}). כאן נדרשת הכרעה מול ממצא \en{Boukrouh} ועמיתיו: הם ממליצים לשלב שיפוט עם חישוב מהנתונים, אבל שיטת האנטרופיה מחשבת משקלות מפיזור הציונים של \textbf{כמה ספקים} באותו קריטריון, ובשלב טרום ההתקשרות מוערך ספק אחד שאין לו נתוני ביצוע. לכן המשקלות נקבעים בשיפוט, ובגלל הרגישות שהם מצאו בשיפוט בלבד, כל משקל נבדק בניתוח רגישות. אם בהמשך יצטבר מאגר הערכות של ספקים רבים, אפשר יהיה לעבור לשיטה המשולבת.

% =====================================================================
\section{הערכת איכות וביצועים של ספקים}\label{sec:lit-quality}

\en{Dey}, \en{Bhattacharya} ו-\en{Ho} בנו מסגרת להערכת ספקים אסטרטגיים, שמשלבת פריסת פונקציית האיכות (\en{QFD}) ו-\en{AHP}, ויישמו אותה בארגון ייצור בבריטניה \cite{dey2015}. לפיהם רוב המחקרים מעריכים ספקים לפי \textbf{מדדי תוצאה} (איכות, עמידה בלוח זמנים, עלות), ויש להוסיף \textbf{מדדים מקדימים}: ניהול הארגון, ניהול סיכונים ופרקטיקות סביבתיות וחברתיות. מדדי תוצאה דורשים היסטוריית עבודה, שאין בשלב טרום ההתקשרות, ולכן איכות מוערכת כאן לפי \textbf{אותות ציבוריים} שקרובים לרעיון המדדים המקדימים.

אות מרכזי הוא הסמכת מערכת ניהול האיכות, כמו \en{ISO 9001} או \en{AS9100} בענף התעופה. \en{Ayyadi} ו-\en{Oulhadj} סקרו 94 מחקרים מהשנים 2000 עד 2023 ומצאו שהשפעת הסמכת \en{ISO 9001} על ביצועי חברות ייצור \textbf{אינה אחידה}, ותלויה בסיבה לבחירה בהסמכה ובאופן השימוש במערכת \cite{ayyadi2024}.

\textbf{ביקורת.} הסמכה מראה שיש מערכת לניהול איכות, לא שהיא עובדת היטב. לכן ממד האיכות משלב כמה אותות: הסמכות בתוקף, לקוחות גדולים וחוזים ממשלתיים, פרסי ספקים, ואותות הפוכים כמו ריקולים והוראות תקינות של רשויות תעופה. התיקוף בודק אם ספקים מוכחים אכן מקבלים בממד זה ציון גבוה יותר.

% =====================================================================
\section{סיכון גיאופוליטי ורמת המדינה}\label{sec:lit-geo}

\en{Ho} ועמיתיו סקרו את ספרות ניהול סיכוני שרשרת האספקה בשנים 2003 עד 2013, ומהסקירה עולה שסיכון לספק נובע גם מגורמים חיצוניים, כמו המדינה שבה הוא פועל \cite{ho2015}. \en{Caldara} ו-\en{Iacoviello} בנו מדד סיכון גיאופוליטי (\en{GPR}) מספירת כתבות על מתחים, מלחמות וטרור \cite{caldara2022}, ובכך הראו שאפשר למדוד סיכון כזה כמותית מתוך טקסט חדשותי, אותו סוג מקור שהסוכנים יעבדו.

שני מחקרים שפיטים מראים שצריך לבדוק גם את \textbf{היחסים} בין מדינת הספק למדינת הלקוח. \en{Bai} ועמיתיו מצאו שלסיכון גיאופוליטי במדינת הספק יש השפעה כפולה על בסיס הספקים של חברות רב-לאומיות, והכיוון שלה תלוי בזיקה הפוליטית בין מדינת החברה למדינת הספק \cite{bai2025}. \en{Charpin} ו-\en{Cousineau} בדקו 2,858 חברות אמריקאיות שרוכשות מ-99 מדינות בשנים 2003 עד 2019, ומצאו שחברות מצמצמות ספקים במדינות שבהן המתיחות עולה \cite{charpin2025}. הממצאים תומכים בקטגוריה \textbf{יחסי המדינה עם ישראל}.

\textbf{ביקורת.} שני המחקרים בוחנים את בסיס הספקים של חברות גדולות לאורך שנים, ולא החלטה על ספק בודד. הם מצדיקים את הקטגוריה, אבל לא קובעים לה משקל; המשקל נקבע ב-\en{AHP}.

% =====================================================================
\section{חיזוי כשל פיננסי של ספקים}\label{sec:lit-fin}

עבודת היסוד היא מודל ה-\en{Z} של \en{Altman}, שמשקלל חמישה יחסים פיננסיים לציון אחד המבחין בין חברות שפשטו רגל לחברות יציבות \cite{altman1968}. אחד היחסים משתמש ב\textbf{שווי השוק} של ההון, ולכן אפשר לחשב אותו ישירות רק לחברה נסחרת. \en{Zhao}, \en{Ouenniche} ו-\en{De Smedt} סקרו את השיטות הסטטיסטיות ושיטות הבינה המלאכותית לחיזוי פשיטת רגל \cite{zhaoj2024}, והסקירה משמשת כאן כמפת השיטות האפשריות לממד הפיננסי.

\textbf{ביקורת: החברות הפרטיות.} כל השיטות האלה זקוקות לנתונים כספיים, ואילו לפי \en{RapidRatings} 22.1\% מהספקים הפרטיים בענף נמצאים בסיכון פיננסי גבוה, לעומת 15.9\% מהציבוריים \cite{rapidratings2026}. הסיכון גבוה דווקא במקום שבו הנתונים חסרים. לכן בחברה פרטית משמשים מדדים עקיפים, כמו אירועים שליליים, שינויי בעלות והסמכות שבוטלו, והדוח מציין כמה מהמידע נמצא.

% =====================================================================
\section{זיהוי אירועים שליליים מטקסט}\label{sec:lit-events}

סינון \textbf{חדשות שליליות} (\en{adverse media}) הוא רכיב קבוע בבדיקות "הכר את הלקוח" (\en{KYC}) ובמניעת הלבנת הון, ושם הוא נעשה בדרך כלל בחיפוש מילות מפתח שמייצר התרעות שווא רבות \cite{chernakov2025}. \en{Chernakov} ועמיתיו בנו סוכן מבוסס מודל שפה (\en{LLM}) עם אחזור מקורות (\en{RAG}), שמחפש ברשת ומחשב לכל נבדק מדד חדשות שליליות, ובדקו אותו על אנשי ציבור, על אנשים מרשימות פיקוח ומרשימות סנקציות ועל שמות "נקיים" \cite{chernakov2025}; זו טיוטה מוקדמת (\en{preprint}), והיא עוסקת באנשים ולא בספקים. \en{Zhang} ועמיתיו חילצו בעזרת מודל שפה 2,843 אירועי שיבוש מ-3,125 כתבות על שרשרת האספקה של רכבים חשמליים \cite{zhang2026}.

\textbf{ביקורת.} המחקרים עוסקים באנשים או בענף שלם, ואף אחד מהם אינו מתרגם אירועים להמלצה על \textbf{ספק בודד}. רובם משייכים כל אירוע לרשימה קבועה של קטגוריות (\textbf{טקסונומיה סגורה}), ולכן אירוע מסוג חדש נדחס לקטגוריה לא מתאימה. בכלי המוצע יש קטגוריה "אחר / חדש" שאדם מאשר, והחומרה נקבעת לפי הפגיעה ביכולת לספק ולא לפי סוג האירוע: מתקפת סייבר שעצרה ייצור היא כשל אספקה, ומתקפה שלא עצרה אותו היא אות סיכון.

% =====================================================================
\section{אבטחת סייבר של ספקים}\label{sec:lit-cyber}

\en{ENISA} מצאה שב-66\% מ-24 מתקפות על שרשראות אספקה הספקים לא ידעו או לא דיווחו איך נפרצו \cite{enisa2021}, ובענף התעופה מתקפת כופרה (\en{ransomware}) על \en{ASCO} עצרה ייצור בכמה מדינות \cite{certeu2019}. המסגרת הרשמית בארה"ב היא \en{NIST SP 800-161 Rev.~1} \cite{nist2022}, ומשרד ההגנה האמריקאי מחייב את ספקיו בהסמכת \en{CMMC} \cite{cmmc2024}, שהיא גם \textbf{אות ציבורי} שאפשר לבדוק בלי לפנות לספק. \en{Gani} ועמיתיו מצאו, בסקר של 130 יצרנים במלזיה, ש\textbf{נראוּת} של שרשרת האספקה קשורה לביצועי אבטחת הסייבר שלה \cite{gani2023}. \en{Slapni{\latinfallback č}ar} ועמיתיו ראיינו 33 מומחים וסקרו 53 נוספים, ובנו תיאוריה של תהליך הערכת סיכוני הסייבר של ספקים \cite{slapnicar2025}.

\textbf{ביקורת.} רוב התהליכים נשענים על \textbf{שאלונים שהספק ממלא בעצמו} ועל גישה למערכות שלו, ואין גישה כזו לפני ההתקשרות; ממצא \en{ENISA} מראה שגם ספק שמשיב לא תמיד יודע מה קרה אצלו. לכן הכלי בודק סייבר \textbf{מבחוץ ובאופן פסיבי בלבד}: הסמכות ודיווחים פומביים על פריצות.

% =====================================================================
\section{מודלי שפה ומערכות מרובות סוכנים בשרשרת אספקה}\label{sec:lit-agents}

\en{Jackson} ועמיתיו הציעו מסגרת רעיונית, שלא נבדקה על נתונים, לשימוש בבינה מלאכותית יוצרת ב-13 תחומי החלטה בשרשרת אספקה \cite{jackson2024}. \en{AlMahri}, \en{Xu} ו-\en{Brintrup} בנו מערכת של \textbf{שבעה סוכנים} שמזהים אותות שיבוש בחדשות, ממפים אותם לרשת ספקים וממליצים על צעדים \cite{almahri2026}. לפיהם המערכת הגיעה ל-\en{F1} של 0.962 עד 0.991, וניתוח ארך 3.83 דקות ועלה כ-0.08 דולר. זו טיוטה שלא עברה שיפוט, שנבדקה על 30 תרחישים \textbf{מסונתזים} בתעשיית הרכב, והיא \textbf{מנטרת ספקים קיימים} ולא בודקת ספק חדש.

\textbf{הסיכון של מודלי שפה.} \en{Huang} ועמיתיו סקרו את תופעת ה\textbf{הזיה} (\en{hallucination}): טקסט שנשמע אמין אבל אינו נכון, שגם אחזור מקורות אינו מבטל לגמרי \cite{huang2025}. בהמלצה על ספק ביטחוני טענה מומצאת היא כשל חמור, ולכן כל טענה מקושרת למקור (\en{REQ-NFR-01}), וטענה בלי מקור אינה נכנסת לציון.

% =====================================================================
\section{מטריצת הספרות}\label{sec:lit-matrix}

טבלה~\ref{tab:litmatrix} מרכזת את המחקרים העיקריים מול הנושאים שהבעיה דורשת.

{\small
\begin{longtable}{|>{\raggedright\arraybackslash}p{2.4cm}|>{\centering\arraybackslash}p{1.12cm}|>{\centering\arraybackslash}p{1.12cm}|>{\centering\arraybackslash}p{1.12cm}|>{\centering\arraybackslash}p{1.12cm}|>{\centering\arraybackslash}p{1.12cm}|>{\centering\arraybackslash}p{1.12cm}|>{\centering\arraybackslash}p{1.12cm}|>{\centering\arraybackslash}p{1.12cm}|}
\caption{מטריצת הספרות: השורות הן מחקרים, והעמודות הן נושאים שהבעיה דורשת}\label{tab:litmatrix}\\
\hline
\textbf{מחקר} & \textbf{טקסט חופשי ומודל שפה} & \textbf{המלצה על ספק בודד} & \textbf{שלב טרום התקש\-רות} & \textbf{איכות ואמינות הספק} & \textbf{חברות פרטיות וכיסוי מידע} & \textbf{סייבר} & \textbf{גיאו\-פוליטי} & \textbf{תיקוף על כשלים אמיתיים} \\ \hline
\endfirsthead
\hline
\textbf{מחקר} & \textbf{טקסט חופשי ומודל שפה} & \textbf{המלצה על ספק בודד} & \textbf{שלב טרום התקש\-רות} & \textbf{איכות ואמינות הספק} & \textbf{חברות פרטיות וכיסוי מידע} & \textbf{סייבר} & \textbf{גיאו\-פוליטי} & \textbf{תיקוף על כשלים אמיתיים} \\ \hline
\endhead
\hline
\endlastfoot
\en{Ho} ועמיתיו \cite{ho2010} & --- & נבדק & נבדק & חלקי & --- & --- & --- & --- \\ \hline
\en{Dey} ועמיתיו \cite{dey2015} & --- & נבדק & --- & נבדק & --- & --- & --- & --- \\ \hline
\en{Altman} \cite{altman1968} & --- & חלקי & --- & --- & --- & --- & --- & נבדק \\ \hline
\en{Bai} ועמיתיו \cite{bai2025} & --- & --- & --- & --- & --- & --- & נבדק & --- \\ \hline
\en{Slapni{\latinfallback č}ar} ועמיתיו \cite{slapnicar2025} & --- & חלקי & --- & --- & --- & נבדק & --- & --- \\ \hline
\en{Chernakov} ועמיתיו \cite{chernakov2025} & נבדק & חלקי & חלקי & --- & --- & --- & --- & חלקי \\ \hline
\en{AlMahri} ועמיתיו \cite{almahri2026} & נבדק & --- & --- & --- & --- & --- & חלקי & --- \\ \hline
\textbf{מאמר נלווה:} \en{Zhao} ועמיתיו \cite{zhao2025} & נבדק & --- & --- & --- & --- & --- & חלקי & --- \\ \hline
\textbf{מאמר הדגל:} \en{Sadeghi} ועמיתיו \cite{sadeghi2026} & חלקי & --- & נבדק & נבדק & --- & --- & חלקי & --- \\ \hline
\textbf{הפרויקט המוצע} & ייבדק & ייבדק & ייבדק & ייבדק & \textbf{ייבדק: זה הפער} & ייבדק & ייבדק & ייבדק \\
\end{longtable}
}

מחקרי \en{MCDM} ואיכות ספקים מגיעים להמלצה ומעריכים חוזקות, אבל מניחים מידע מלא והיסטוריית ביצוע. מחקרי מודלי השפה מחפשים רק סיכונים. מאמר הדגל מחבר בין העולמות, אבל אינו אוסף מידע ציבורי ואינו מטפל במידע חסר. אף מחקר אינו מדווח כמה מהמידע על הספק נמצא.

% =====================================================================
\section{מאמר הדגל}\label{sec:flagship}

\textbf{הבחירה.} מאמר הדגל הוא המאמר של \en{Sadeghi}, \en{Dehghani} ו-\en{Kabadayi} (2026) בכתב העת \en{Annals of Operations Research} \cite{sadeghi2026}. הוא \textbf{עוסק באותה בעיה}: בחירת ספקים לפי כמה קריטריונים בעזרת מודל שפה ושקלול ב-\en{AHP}, כשאיכות ואמינות הם אחד הממדים; הוא \textbf{עדכני ופורסם בכתב עת שפיט מוביל}; ו\textbf{אפשר לשחזר אותו}, כי השלבים, משוואות ה-\en{AHP} וסף העקביות מתוארים במלואם, והתוצאות הושוו לשיפוט מומחים במבחן סטטיסטי.

\textbf{המודל.} המסגרת \en{AIDM} (\en{AI-driven Decision-Making}) מחליפה במודל שפה את פאנל המומחים שממלא את ה-\en{AHP}, בארבעה שלבים: (1) המודל יוצר חמישה מומחים וירטואליים עם רקע שונה; (2) כל מומחה מציע קריטריונים, והחשובים ביותר מקובצים להיררכיה; (3) המומחים משווים בזוגות, והשוואות לא עקביות חוזרות עד חמש פעמים; (4) המשקלות מושווים למשקלות שקבעו מומחי רכש אמיתיים. ההיררכיה שהתקבלה כוללת שלושה ממדים: \textbf{איכות} (איכות מוצר, אמינות, תקשורת, גמישות ומשוב לקוחות), \textbf{עלות} (מחיר, זמן אספקה, עדכוני מחיר ותנאי חוזה) ו\textbf{מאפייני הספק} (השפעה סביבתית, נגישות, \textbf{מדינת המקור}, ממליצים, מערכות ניהול וניסיון גלובלי).

\textbf{המשוואות השולטות.} כל מעריך ממלא מטריצת השוואות $A$ בגודל $n\times n$, שבה $a_{ij}$ אומר כמה קריטריון $i$ חשוב מקריטריון $j$. המשקלות הם הווקטור העצמי של הערך העצמי הגדול ביותר $\lambda_{\max}$, והעקביות נמדדת ביחס $\mathrm{CR}$, כש-$\mathrm{RI}$ הוא ערך ייחוס שתלוי בגודל המטריצה \cite{saaty1990}:
\begin{equation}\label{eq:ahp}
A\,\mathbf{w}=\lambda_{\max}\,\mathbf{w},\qquad \sum_{i=1}^{n} w_i = 1
\end{equation}
\begin{equation}\label{eq:cr}
\mathrm{CI}=\frac{\lambda_{\max}-n}{n-1},\qquad
\mathrm{CR}=\frac{\mathrm{CI}}{\mathrm{RI}}\le 0.10
\end{equation}
סף העקביות במאמר הוא 0.10, והוא מאומץ גם כאן. כדי להשוות בין משקלות המומחים האנושיים $w^{H}$ לווירטואליים $w^{AI}$, חושב לכל קריטריון $|w^{H}_i-w^{AI}_i|$, וההבדל נבדק במבחן \en{Wilcoxon} לזוגות \cite{sadeghi2026}.

\textbf{הניסוי והתוצאות.} המקרה הוא בחירת ספקים בחברה לייצור מכונות כבדות; המומחים הווירטואליים נבנו עם \en{\mbox{ChatGPT-4}}, וקבוצת ההשוואה כללה 14 אנשי רכש (8 ברכש המקומי ו-6 בבינלאומי) \cite{sadeghi2026}. שתי הקבוצות עמדו בסף העקביות, והאנושיים היו עקביים יותר ($\mathrm{CR}=0.030$ ו-$0.060$); 85\% מהמומחים הווירטואליים הגיעו ליחס תקין. מבחן \en{Wilcoxon} לא מצא הבדל מובהק ($\alpha=0.05$) בין קבוצות המשקלות, אבל ברמת הקריטריון היו הבדלים עקביים (טבלה~\ref{tab:sadeghi-results}).

{\small
\begin{longtable}{|>{\raggedright\arraybackslash}p{3.6cm}|c|c|c|c|}
\caption{דוגמאות למשקלות במאמר הדגל: מומחים אנושיים מול וירטואליים \cite{sadeghi2026}}\label{tab:sadeghi-results}\\
\hline
\textbf{קריטריון} & \textbf{מקומי: אנושי} & \textbf{מקומי: וירטואלי} & \textbf{בינלאומי: אנושי} & \textbf{בינלאומי: וירטואלי} \\ \hline
\endfirsthead
\hline
\textbf{קריטריון} & \textbf{מקומי: אנושי} & \textbf{מקומי: וירטואלי} & \textbf{בינלאומי: אנושי} & \textbf{בינלאומי: וירטואלי} \\ \hline
\endhead
אמינות & 0.114 & 0.064 & 0.139 & 0.075 \\ \hline
עדכוני מחיר & 0.099 & 0.031 & 0.107 & 0.033 \\ \hline
מדינת המקור & 0.029 & 0.065 & 0.024 & 0.066 \\ \hline
מודעות סביבתית & 0.024 & 0.050 & 0.016 & 0.051 \\* \hline
\end{longtable}
}

אנשי הרכש נתנו משקל גבוה פי שניים בערך ל\textbf{אמינות}, ומודל השפה משקל גבוה פי שניים בערך ל\textbf{מדינת המקור}. אף שהמבחן לא מצא הבדל ברמת כלל המשקלות, ההבדלים חשובים דווקא בגלל הקריטריונים שבהם הופיעו: אלה שני הנושאים הרגישים ביותר לתעשייה ביטחונית. לכן במתודולוגיה המוצעת המשקלות נקבעים בשיפוט אנושי המנומק בספרות, ולא על ידי הסוכנים. מהמאמר נלקחים המבנה ההיררכי, משוואות ה-\en{AHP}, סף העקביות וממדי האיכות והאמינות; הרעיון של מודל שפה כמעריך משקלות נדחה.

\textbf{המגבלות.} המחברים מציינים שמודל השפה עלול להמציא קריטריונים, ושאין ערובה שהרשימה מלאה \cite{sadeghi2026}. בנוסף זוהו שש מגבלות: (1) המאמר \textbf{עוצר במשקלות} ואינו מעריך אף ספק בפועל; (2) המסגרת \textbf{אינה אוספת מידע} על הספק; (3) \textbf{אין טיפול במידע חסר}; (4) \textbf{סיכון הוא משקל בלבד}, כך שספק ממדינה מסוכנת יכול "לפצות" במחיר נמוך, ואין שער חסימה, אירועים חריגים או סייבר; (5) התיקוף הוא \textbf{מול מומחים ולא מול המציאות}; (6) \textbf{חברה אחת וענף אחד}.

% =====================================================================
\section{המאמר הנלווה: זיהוי סיכונים מחדשות}\label{sec:companion}

מאמר הדגל עוסק בעיקר בפרופיל החוזקות ובמשקלות. לחלק דגלי הסיכון משמש מאמר נלווה: \en{Zhao} ועמיתיו (2025), \en{Service Oriented Computing and Applications} \cite{zhao2025}, שבודק אם מודלי שפה מזהים נכון סיכוני ספקים מתוך כתבות אמיתיות, כלומר את משימת סוכני הסיכון.

המסגרת \en{LARD-SC} אוספת כתבות על כל ספק מ-\en{Google News}; מודל שפה מסכם כל כתבה, מזהה סיכון ומעריך את סבירותו והשפעתו; וכל סיכון ממופה לטקסונומיה של \en{Cambridge} לסיכונים עסקיים. מתוך 676 כתבות על 188 ספקים בשכבה הראשונה של \en{Apple} (2020--2024) נדגמו 120, חמישה מודלים קיבלו אותן עם אותו פרומפט, ושני זוגות מומחים בדקו כל סיכון. חושבו שיעור הסיכונים שאושרו (\en{RVR}), שאושרו חלקית (\en{PRR}), הזיהויים השגויים (\en{FIR}) ומדד ביצוע יחסי (\en{RPI}) (טבלה~\ref{tab:zhao-results}).

{\small
\begin{longtable}{|l|c|c|c|c|}
\caption{תוצאות המאמר הנלווה: ממוצע שני זוגות המומחים \cite{zhao2025}}\label{tab:zhao-results}\\
\hline
\textbf{מודל} & \en{RVR} & \en{PRR} & \en{FIR} & \en{RPI} \\ \hline
\endfirsthead
\hline
\textbf{מודל} & \en{RVR} & \en{PRR} & \en{FIR} & \en{RPI} \\ \hline
\endhead
\hline
\endlastfoot
\en{GPT-4o} & 0.8875 & 0.0917 & 0.0208 & 0.8484 \\ \hline
\en{GPT-4o mini} & 0.8750 & 0.1042 & 0.0208 & 0.8334 \\ \hline
\en{Claude 3.5 Sonnet} & 0.8708 & 0.1042 & 0.0250 & 0.8250 \\ \hline
\en{Claude 3.5 Haiku} & 0.8250 & 0.1542 & 0.0208 & 0.7734 \\ \hline
\en{GPT-3.5 turbo} & 0.4625 & 0.2417 & 0.2958 & 0.1184 \\
\end{longtable}
}

במאמר מוגדר \en{$\mathrm{RPI}=w_1\mathrm{RVR}-w_2\mathrm{PRR}-w_3\mathrm{FIR}$} עם המשקלות \en{$w_1=1,\ w_2=0.5,\ w_3=1$}, אבל ערכי ה-\en{RPI} בטבלה 3 של המאמר, שמוצגים כאן, מתקבלים רק עם \en{$w_2=0.2$}; למשל ל-\en{GPT-4o}: \en{$0.8875-0.2\cdot0.0917-0.0208=0.8484$}. דירוג המודלים זהה בשתי הדרכים. ארבעת המודלים החדשים הגיעו לשיעור אימות של 0.82 עד 0.89 עם פחות מ-3\% זיהויים שגויים, והישן נכשל בכמעט 30\% מהמקרים: \textbf{בחירת המודל משנה את התוצאה מהותית}, ולכן התיקוף ייעשה בשני מודלים לפחות (\en{T4.1}, פרק~\ref{ch:workplan}).

\textbf{המגבלות.} מעבר להשפעת ניסוח הפרומפט ולתלות בזמינות המידע, שהמחברים מציינים \cite{zhao2025}: (1) \textbf{רק דיוק, בלי כיסוי}: אין מדד לסיכונים שהוחמצו, ובספק ביטחוני החמצה חמורה מהתרעת שווא; (2) המומחים בדקו אם הסיכון \textbf{סביר}, ולא אם התממש; (3) הספקים \textbf{גלויים}, וכשאין כתבות, היעדר חדשות רעות עלול להתפרש כהיעדר סיכון; (4) הניתוח ברמת \textbf{הכתבה}, בלי ציון ספק ובלי שערי חסימה; (5) \textbf{רק סיכונים}, בלי חוזקות.

% =====================================================================
\section{פער הידע}\label{sec:gap}

\begin{quote}
\textbf{עד כה לא נבדק} אם מערכת מבוססת מודלי שפה, שנשענת על מידע ציבורי בלבד, יכולה להפיק \textbf{המלצת התקשרות על ספק בודד} בשלב טרום ההתקשרות, להעריך בה גם את \textbf{איכות הספק ואמינותו} ולא רק סיכונים, ולדווח במפורש \textbf{כמה מהמידע הנדרש נמצא}, כשהספק הוא חברה פרטית שעליה מעט מידע, ולתקף את ההמלצה מול \textbf{מקרי כשל אמיתיים}.
\end{quote}

הפרויקט עונה על המגבלות שתוארו: הוא ממשיך ממשקלות להערכה מלאה של ספק מסוים; הסוכנים אוספים מידע ציבורי וכל טענה מקושרת למקור; סנקציות, בקרת יצוא ומדינה עוינת הם שער פסילה ולא קריטריון בשקלול; ציון הכיסוי מונע מהיעדר כתבות להתפרש כהיעדר סיכון; בחברה פרטית משמשים מדדים עקיפים; והתיקוף נעשה מול מקרי כשל ציבוריים ומול ספקים מוכחים, כך שנמדדות גם החמצות.

% ---------- פרק 4: חקר השוק והמוצרים הקיימים ----------
% 22.11.2026: גוף שלישי, קיצור, איורים בעברית. לפי פרק 6 במדריך.
% מחירי המוצרים לא פורסמו – לא להמציא. הציונים בטבלה 4.2 הם הערכה לפי התיעוד הפומבי, לא מדידה.
\chapter{חקר השוק והמוצרים הקיימים}\label{ch:market}

הפרק עונה על שתי שאלות: כמה גדול השוק לתוכנות לבדיקת ספקים, ואילו מוצרים קיימים ומה הם בודקים, את איכות הספק, את הסיכונים שבו או את שניהם.

\section{היקף השוק: שלוש שכבות}\label{sec:tam}

\textbf{\en{TAM}: השוק העולמי לתוכנות ניהול סיכוני ספקים.} ה-\en{TAM} נמדד לפי שוק התוכנות לניהול סיכוני ספקים וצדדים שלישיים (\en{Third-Party Risk Management}), שבהן חברות בודקות ומנטרות את ספקיהן. לפי \en{WiseGuy Reports} השוק הוערך ב-5.93 מיליארד דולר ב-2025, וצפוי לצמוח ל-12.5 מיליארד עד 2035, בשיעור צמיחה שנתי ממוצע (\en{CAGR}) של 7.7\% \cite{wiseguy2026}. זהו דוח מסחרי שאינו מפרסם את שיטת החישוב, וחברות מחקר אחרות מגדירות את השוק אחרת, ולכן המספר משמש כסדר גודל. אינדיקציה בלתי תלויה לבשלות השוק היא דוח \en{Magic Quadrant} נפרד של \en{Gartner} לפתרונות ניהול סיכוני ספקים, שפורסם במאי 2026 \cite{everstream2026}.

\textbf{\en{SAM}: התעשייה הביטחונית.} הפלח הוא \textbf{יצרניות ביטחוניות ואוויריות}, שחייבות לבדוק סנקציות ובקרת יצוא ומוגבלות במידע שמותר להן להעלות לענן. לפי \en{SIPRI}, הכנסות 100 יצרניות הנשק הגדולות הגיעו ב-2024 לשיא של 679 מיליארד דולר, עלייה של 5.9\%, ותשע מהן במזרח התיכון \cite{sipri2025}. \en{SIPRI} והדוחות הכספיים של היצרניות אינם מפרטים את ההוצאה על כלים לבדיקת ספקים, ולא נמצא לה מקור ציבורי; לכן ה-\en{SAM} מוגדר לפי \textbf{מספר ארגונים}, כ-100 היצרניות הגדולות, ולא לפי סכום כספי.

\textbf{\en{SOM}: ארגון אחד בשנה הראשונה.} הנתח שאפשר לתפוס בפועל הוא יחידות הרכש והפיתוח העסקי של תע"א, שעבדה בסוף 2022 עם כ-5,100 ספקים \cite{mevaker2024}. הערך השנתי $V$ מחושב מהחיסכון בשעות עבודה:
\begin{equation}\label{eq:som}
V = N \cdot \left(t_{\mathrm{manual}} - t_{\mathrm{tool}}\right) \cdot c_{h} - C_{\mathrm{run}}
\end{equation}
כאשר $N$ הוא מספר הספקים החדשים שנבדקים בשנה, $t_{\mathrm{manual}}$ ו-$t_{\mathrm{tool}}$ שעות הבדיקה היום ועם הכלי, $c_{h}$ עלות שעת עבודה ו-$C_{\mathrm{run}}$ עלות ההפעלה השנתית. $N$, $t_{\mathrm{manual}}$ ו-$c_{h}$ הם נתונים פנימיים של הארגון שלא פורסמו, ו-$t_{\mathrm{tool}}$ ו-$C_{\mathrm{run}}$ נמדדים בהרצת הכלי על מקרי התיקוף (פרק~\ref{ch:design}). אינדיקציה גסה לעלות ההפעלה: במערכת של \en{AlMahri} ועמיתיו ניתוח מלא עלה כ-0.08 דולר וארך כ-4 דקות \cite{almahri2026}, כלומר עלות זניחה ביחס לשעת אנליסט.

\section{סקירת מוצרים}\label{sec:products}

נסקרו ארבעה מוצרים מסחריים, שכל אחד מייצג גישה אחרת. התיאורים מבוססים על מה שהחברות מפרסמות; המוצרים לא נבדקו בפועל, ואף חברה אינה מפרסמת מחירון.
\begin{itemize}
\item \textbf{\en{RapidRatings} (פיננסי):} ציון בריאות פיננסית (\en{FHR}) בסולם 0 עד 100 מדוחות כספיים, שחוזה חדלות פירעון ב-12 החודשים הבאים. לדירוג חברה פרטית, החברה צריכה \textbf{למסור את דוחותיה} למערכת של \en{RapidRatings} \cite{rrfaq}, כלומר נדרש שיתוף פעולה של הספק.
\item \textbf{\en{Interos} (רב-ממדי):} ציון מאוחד (\en{i-Score}) שמכסה סייבר, רגולציה, גיאופוליטיקה, אסונות, פיננסים ו-\en{ESG}, על יותר מ-250 מיליון ישויות \cite{interosgov}. ב-2023 נבחר לספק לצי האמריקאי ניהול סיכוני שרשרת אספקה \cite{interosnavy2023}, ולכן הוא הקרוב ביותר לשוק הביטחוני.
\item \textbf{\en{Bitsight} (סייבר):} דירוג אבטחה בסולם 250 עד 900 ממידע שנאסף מבחוץ, \textbf{בלי מידע מהחברה הנבדקת} \cite{bitsight2026}, בדומה לגישה הפסיבית של ממד הסייבר המוצע.
\item \textbf{\en{Everstream} ו-\en{Prewave} (ניטור אירועים):} ניתוח מקורות פתוחים בבינה מלאכותית להתרעה על אירועים אצל ספקים. שתיהן מדווחות שדורגו כ-\en{Leaders} ב-\en{Gartner} 2026 \cite{everstream2026,prewave2026}, והן מוצגות כמוצר אחד.
\end{itemize}

טבלה~\ref{tab:products} משווה אותם למסגרת \en{AIDM} של מאמר הדגל, לתהליך הידני ולפתרון המוצע.

{\footnotesize
\setlength{\tabcolsep}{3pt}
\begin{longtable}{|>{\raggedright\arraybackslash}p{1.75cm}|>{\raggedright\arraybackslash}p{1.55cm}|>{\raggedright\arraybackslash}p{1.55cm}|>{\raggedright\arraybackslash}p{1.45cm}|>{\raggedright\arraybackslash}p{1.55cm}|>{\raggedright\arraybackslash}p{1.55cm}|>{\raggedright\arraybackslash}p{1.45cm}|>{\raggedright\arraybackslash}p{1.65cm}|}
\caption{השוואה רב-פרמטרית בין המוצרים, מאמר הדגל, התהליך הידני והפתרון המוצע}\label{tab:products}\\
\hline
\textbf{פרמטר} & \en{RapidRatings} & \en{Interos} & \en{Bitsight} & \en{Everstream} / \en{Prewave} & \textbf{מאמר הדגל} (\en{AIDM}) & \textbf{תהליך ידני} & \textbf{הפתרון המוצע} \\ \hline
\endfirsthead
\hline
\textbf{פרמטר} & \en{RapidRatings} & \en{Interos} & \en{Bitsight} & \en{Everstream} / \en{Prewave} & \textbf{מאמר הדגל} (\en{AIDM}) & \textbf{תהליך ידני} & \textbf{הפתרון המוצע} \\ \hline
\endhead
עיקרון הפעולה & מודל על דוחות כספיים & ציון מאוחד מכמה תחומי סיכון & סריקה חיצונית של תשתיות & ניתוח טקסט ממקורות פתוחים & \en{AHP} עם מומחים וירטואליים & חיפוש ידני ושיקול דעת & סוכן לכל ממד, שער ושקלול \\ \hline
ממדים מכוסים & פיננסי & רוב הממדים & סייבר & אירועים, חלקית פיננסי וגיאופוליטי & איכות, עלות ומאפייני ספק & תלוי באנליסט & שבעה ממדים, שני חלקים \\ \hline
איכות ואמינות הספק & לא & לא פורסם & לא & לא פורסם & כן, בשקלול & לפי האנליסט & ממד ייעודי \\ \hline
ספק פרטי בלי דוחות & דורש שהספק ימסור & לא פורסם & לא רלוונטי & לפי כיסוי בחדשות & לא אוסף מידע & חלקי & מדדים עקיפים וכיסוי \\ \hline
דיווח על מידע חסר & לא פורסם & לא פורסם & לא פורסם & לא פורסם & לא & לא שיטתי & ציון כיסוי ורשימה ידנית \\ \hline
שלב בתהליך & ניטור ובחירה & ניטור & בחירה וניטור & ניטור & בחירה & טרום התקשרות & \textbf{טרום התקשרות} \\ \hline
מקור לכל טענה & לא פורסם & לא פורסם & לא פורסם & לא פורסם & לא & תלוי באנליסט & חובה \\ \hline
עלות & לא פורסמה & לא פורסמה & לא פורסמה & לא פורסמה & לא פורסמה & שעות אנליסט & קריאות למודל \\ \hline
נקודת תורפה & תלוי בספק הפרטי & סגור, לא ברור איך מגיע לציון & ממד אחד & ספקים קיימים & עוצר במשקלות & איטי, לא אחיד & לא נבדק עדיין \\ \hline
\end{longtable}
}

\textbf{מה הטבלה מראה.} המוצרים חזקים בממד שלהם, ו-\en{Interos} מכסה את רובם ופועל בשוק הביטחוני. אבל בשלוש שורות כמעט כל המשבצות הן "לא" או "לא פורסם": \textbf{איכות ואמינות}, \textbf{דיווח על מידע חסר} ו\textbf{מקור לכל טענה} (\en{REQ-FUN-03} ו-\en{REQ-NFR-01}). המוצרים בנויים לחפש סיכונים, ולא לענות על השאלה "האם הספק טוב". מאחר שהבדיקה נעשתה על התיעוד הפומבי בלבד, הממצא הוא "לא נמצא בתיעוד", ולא "לא קיים".

\section{השוואה גרפית}\label{sec:market-graphs}

שני הגרפים נגזרים מטבלה~\ref{tab:products} דרך טבלת הציונים~\ref{tab:axes}, על \textbf{שישה נושאים}: איכות והתאמה טכנית אוחדו, כי אף מוצר אינו מפריד ביניהן; עמידה ביעדים הושמטה, כי אף מוצר אינו מפרסם שהוא בודק אותה; ודיווח הכיסוי נוסף. כל נושא מקבל 0 (לא מכוסה), 1 (חלקית) או 2 (מכוסה); "לא פורסם" נספר 0 ומסומן בכוכבית. הפתרון המוצע מוצג לפי \textbf{היעד שלו}, ולצידו המסגרת \en{LARD-SC} של המאמר הנלווה. מאמר הדגל אינו מופיע, כי אינו אוסף מידע על ספקים.

{\small
\begin{longtable}{|>{\raggedright\arraybackslash}p{3.2cm}|c|c|c|c|c|c|c|}
\caption{ציוני הכיסוי לגרפים (0 עד 2). כוכבית: לא פורסם, נספר 0}\label{tab:axes}\\
\hline
\textbf{מוצר} & \textbf{פיננסי} & \textbf{סייבר} & \textbf{גיאו\-פוליטי} & \textbf{אירועים} & \textbf{איכות וטכני} & \textbf{דיווח כיסוי} & \textbf{סכום} \\ \hline
\endfirsthead
\hline
\textbf{מוצר} & \textbf{פיננסי} & \textbf{סייבר} & \textbf{גיאו\-פוליטי} & \textbf{אירועים} & \textbf{איכות וטכני} & \textbf{דיווח כיסוי} & \textbf{סכום} \\ \hline
\endhead
\en{RapidRatings} & 2 & 0 & 0 & 0 & 0 & 0* & 2 \\ \hline
\en{Interos} & 2 & 2 & 2 & 1 & 0* & 0* & 7 \\ \hline
\en{Bitsight} & 0 & 2 & 0 & 0 & 0 & 0* & 2 \\ \hline
\en{Everstream} / \en{Prewave} & 1 & 0 & 1 & 2 & 0* & 0* & 4 \\ \hline
\en{LARD-SC} (מאמר נלווה) & 1 & 0 & 1 & 2 & 0 & 0 & 4 \\ \hline
הפתרון המוצע (יעד) & 2 & 2 & 2 & 2 & 2 & 2 & 12 \\ \hline
\end{longtable}
}

\begin{figure}[htbp]
\centering
\begin{tikzpicture}[scale=0.85]
\def\R{3}
\foreach \r in {1,2}{
  \draw[gray!40] (90:\r*\R/2) \foreach \a in {150,210,270,330,30,90}{ -- (\a:\r*\R/2)};
}
\foreach \a/\lab in {90/פיננסי,150/סייבר,210/גיאופוליטי,270/אירועים חריגים,330/איכות והתאמה טכנית,30/דיווח כיסוי}{
  \draw[gray!60] (0,0) -- (\a:\R);
  \node[font=\small] at (\a:\R+1.25) {\foreignlanguage{hebrew}{\lab}};
}
% Interos: 2,2,2,1,0,0
\draw[blue,thick,fill=blue,fill opacity=0.08] (90:\R) -- (150:\R) -- (210:\R) -- (270:\R/2) -- (330:0) -- (30:0) -- cycle;
% RapidRatings: 2,0,0,0,0,0
\draw[orange,thick,dashed] (90:\R) -- (150:0) -- (210:0) -- (270:0) -- (330:0) -- (30:0) -- cycle;
% LARD-SC: 1,0,1,2,0,0
\draw[violet,thick,dotted] (90:\R/2) -- (150:0) -- (210:\R/2) -- (270:\R) -- (330:0) -- (30:0) -- cycle;
% הפתרון המוצע (יעד): 2 בכולם
\draw[red!70!black,very thick] (90:\R) \foreach \a in {150,210,270,330,30}{ -- (\a:\R)} -- cycle;
\end{tikzpicture}

\vspace{1mm}
{\small
\textcolor{red!70!black}{\rule[2pt]{6mm}{1.5pt}} הפתרון המוצע (יעד) \quad
\textcolor{violet}{\rule[2pt]{6mm}{1pt}} \en{LARD-SC} \quad
\textcolor{orange}{\rule[2pt]{6mm}{1pt}} \en{RapidRatings} \quad
\textcolor{blue}{\rule[2pt]{6mm}{1pt}} \en{Interos}}
\caption{תרשים רדאר: כיסוי של שישה נושאים (0 עד 2), לפי טבלה~\ref{tab:axes}}\label{fig:radar}
\end{figure}

איור~\ref{fig:radar} מראה ש-\en{Interos} מכסה את רוב הנושאים, וההבדל בינו לבין הפתרון המוצע הוא בשני צירים: \textbf{איכות הספק והתאמתו הטכנית}, ו\textbf{דיווח על כיסוי המידע}.

\begin{figure}[htbp]
\centering
\begin{tikzpicture}
\begin{axis}[
  width=10.5cm, height=6.3cm,
  xmin=0, xmax=13, ymin=-0.5, ymax=2.6,
  xlabel={\foreignlanguage{hebrew}{רוחב הכיסוי: סכום הציונים בטבלה \ref{tab:axes}}},
  ylabel={\foreignlanguage{hebrew}{גישה למידע}},
  ytick={0,1,2},
  yticklabels={\foreignlanguage{hebrew}{דורש מידע מהספק},\foreignlanguage{hebrew}{מידע מסחרי בתשלום},\foreignlanguage{hebrew}{מידע ציבורי בלבד}},
  yticklabel style={font=\footnotesize},
  grid=both, grid style={gray!20},
]
\addplot[only marks, mark=*, mark size=3pt, blue] coordinates {(7,1)};
\node[font=\footnotesize, anchor=south west] at (axis cs:7.1,1.08) {Interos};
\addplot[only marks, mark=*, mark size=3pt, orange] coordinates {(2,0)};
\node[font=\footnotesize, anchor=south] at (axis cs:2,0.08) {RapidRatings};
\addplot[only marks, mark=*, mark size=3pt, teal] coordinates {(2,1)};
\node[font=\footnotesize, anchor=north] at (axis cs:2,0.92) {Bitsight};
\addplot[only marks, mark=*, mark size=3pt, gray] coordinates {(4,1)};
\node[font=\footnotesize, anchor=south] at (axis cs:4,1.08) {Everstream/Prewave};
\addplot[only marks, mark=square*, mark size=3pt, violet] coordinates {(4,2)};
\node[font=\footnotesize, anchor=south] at (axis cs:4,2.08) {LARD-SC};
\addplot[only marks, mark=o, mark size=4pt, very thick, red!70!black] coordinates {(12,2)};
\node[font=\footnotesize, anchor=south] at (axis cs:11.4,2.08) {\foreignlanguage{hebrew}{הפתרון המוצע}};
\end{axis}
\end{tikzpicture}
\caption{גרף פיזור: רוחב הכיסוי מול סוג הגישה למידע. הפתרון המוצע (עיגול ריק) מוצג לפי היעד שלו}\label{fig:scatter}
\end{figure}

איור~\ref{fig:scatter} מראה שבדרך כלל כיסוי רחב דורש מידע מסחרי בתשלום, ודירוג פיננסי של חברה פרטית דורש שהספק ימסור דוחות. הפתרון המוצע מכוון לפינה הימנית העליונה, \textbf{כיסוי רחב מתוך מידע ציבורי בלבד}. המחיר הוא שהכיסוי לא יהיה מלא בכל מקרה, ולכן ציון הכיסוי הוא חלק מהפלט.

\textbf{מסקנה.} הכלי המוצע אינו מתחרה ב-\en{Interos} על היקף הנתונים. הוא מוסיף שכבה שלא נמצאה בתיעוד של אף מוצר: \textbf{הערכה מלאה}, חוזקות לצד סיכונים, עם \textbf{המלצה שקופה} שמציגה מקור לכל טענה ומדווחת כמה מהמידע נמצא. הוא אינו דורש נתונים פנימיים, אך באב-הטיפוס שם הספק והרכיב נשלחים למודל שפה חיצוני, וזה עצמו מידע על כוונות רכש; לכן אבטחת מידע היא קריטריון בבחירת החלופה (פרק~\ref{ch:design}) וסיכון בפרק~\ref{ch:risks}.

% ---------- פרק 5: חלופות תכן, בחירת פתרון וחדשנות ----------
% חישוב המטריצה: א 3.25, ב 2.95, ג 4.25. תרחיש אבטחה 0.30/הסבריות 0.05: א 3.40, ב 2.80, ג 3.95. מפנה: אבטחה ~0.43 (א עוקפת ג), ישימות ~0.45 (ב עוקפת ג).
\chapter{חלופות תכן, בחירת פתרון וחדשנות}\label{ch:design}

הפרק מציג שלוש חלופות שונות בעיקרון, בוחר ביניהן במטריצת החלטה משוקללת, ומפרט את הארכיטקטורה, את מקורות המידע ואת החדשנות של החלופה שנבחרה.

\section{שלוש חלופות שונות בעיקרון}\label{sec:alternatives}

החלופות (טבלה~\ref{tab:alternatives}) נבדלות זו מזו באופן שבו המידע נאסף ומעובד.

{\small
\begin{longtable}{|>{\raggedright\arraybackslash}p{3.6cm}|>{\raggedright\arraybackslash}p{4.8cm}|>{\raggedright\arraybackslash}p{5.1cm}|}
\caption{שלוש חלופות התכן, עם היתרון והמחיר של כל אחת}\label{tab:alternatives}\\
\hline
\textbf{חלופה} & \textbf{היתרון} & \textbf{המחיר} \\ \hline
\endfirsthead
\hline
\textbf{חלופה} & \textbf{היתרון} & \textbf{המחיר} \\ \hline
\endhead
\textbf{א: ניקוד מובנה עם איסוף ידני.} אנליסט ממלא תבנית קבועה, והציון מחושב לפי משקלות \en{AHP} & שקוף, לא תלוי במודל שפה, ולא שולח מידע לשירות חיצוני & איטי ותלוי באנליסט; אין זיהוי פתוח של אירועים ומעקב שיטתי אחר מידע חסר \\ \hline
\textbf{ב: סוכן יחיד עם אחזור מקורות (\en{RAG}).} פרומפט אחד מבקש מהמודל לחפש, לנתח ולהמליץ & מהיר לבנייה וזול להפעלה & קשה לבדוק איך הגיע להמלצה; אין שערים קשיחים, וסיכון ההזיה גבוה \cite{huang2025}; ממד בלי מידע עלול "להיעלם" \\ \hline
\textbf{ג: מערכת מרובת סוכנים עם שערים דטרמיניסטיים.} סוכן לכל ממד, בדיקת סנקציות בקוד ושקלול בנוסחה גלויה & כל ממד מדווח בנפרד על ציון, כיסוי ומקורות; השערים לא תלויים במודל; אפשר להחליף סוכן או מודל בנפרד & מורכבת יותר לבנייה, ויקרה יותר בקריאות למודל \\* \hline
\end{longtable}
}

\section{מטריצת החלטה משוקללת}\label{sec:matrix}

כל קריטריון ומשקלו נגזרו מהדרישות בפרק~\ref{ch:problem} לפני ניקוד החלופות:
\begin{itemize}
\item \textbf{הסבריות ומקור לכל טענה (0.20):} \en{REQ-NFR-01} ו-\en{REQ-FUN-09}.
\item \textbf{טיפול במידע חסר (0.20):} \en{REQ-FUN-03} עד \en{REQ-FUN-05}, לב פער הידע.
\item \textbf{הערכת חוזקות וזיהוי סיכונים (0.20):} \en{REQ-PER-01}, \en{REQ-FUN-06} ו-\en{REQ-FUN-10}.
\item \textbf{אבטחת מידע (0.15):} \en{REQ-CON-01} ו-\en{REQ-CON-03}; כמה מידע על כוונות הרכש יוצא לשירות חיצוני, והאם אפשר להריץ על מודל מקומי.
\item \textbf{ישימות (0.10):} אב-טיפוס עובד בסוף שבוע 10.
\item \textbf{גמישות להתאמה בתעשייה האווירית (0.10):} כמה קל להחליף סוכן, מקור או מודל בלי לשנות את שאר המערכת.
\item \textbf{עלות לדוח (0.05):} משקל נמוך, כי עלות קריאות למודל זניחה ביחס לשעת אנליסט (משוואה~\ref{eq:som}).
\end{itemize}

כל חלופה מקבלת ציון $c_{ij}$ בין 1 ל-5 בכל קריטריון, והציון הכולל של חלופה $j$ הוא:
\begin{equation}\label{eq:pugh}
S_j=\sum_{i=1}^{n} w_i\, c_{ij},\qquad \sum_{i=1}^{n} w_i = 1
\end{equation}

{\small
\begin{longtable}{|>{\raggedright\arraybackslash}p{4.6cm}|c|c|c|c|}
\caption{מטריצת החלטה משוקללת: ציון, ובסוגריים הציון המשוקלל}\label{tab:pugh}\\
\hline
\textbf{קריטריון} & \textbf{משקל} & \textbf{א: ידני} & \textbf{ב: סוכן יחיד} & \textbf{ג: רב-סוכני} \\ \hline
\endfirsthead
\hline
\textbf{קריטריון} & \textbf{משקל} & \textbf{א: ידני} & \textbf{ב: סוכן יחיד} & \textbf{ג: רב-סוכני} \\ \hline
\endhead
הסבריות ומקור לכל טענה & 0.20 & 4 (0.80) & 3 (0.60) & 5 (1.00) \\ \hline
טיפול במידע חסר & 0.20 & 2 (0.40) & 2 (0.40) & 5 (1.00) \\ \hline
הערכת חוזקות וסיכונים & 0.20 & 3 (0.60) & 3 (0.60) & 4 (0.80) \\ \hline
אבטחת מידע & 0.15 & 5 (0.75) & 2 (0.30) & 3 (0.45) \\ \hline
ישימות & 0.10 & 4 (0.40) & 5 (0.50) & 3 (0.30) \\ \hline
גמישות להתאמה & 0.10 & 2 (0.20) & 3 (0.30) & 5 (0.50) \\ \hline
עלות לדוח & 0.05 & 2 (0.10) & 5 (0.25) & 4 (0.20) \\ \hline
\textbf{ציון סופי} & \textbf{1.00} & \textbf{3.25} & \textbf{2.95} & \textbf{4.25} \\* \hline
\end{longtable}
}

\textbf{ביסוס הציונים.}
\begin{itemize}
\item \textbf{הסבריות:} ג' מקבלת 5, כי כל סוכן מחזיר מקור לכל טענה והשערים בקוד; ב' 3, כי מודל יחיד עלול להפיק טענות בלי מקור \cite{huang2025}; א' 4, כי אנליסט מתעד מקורות, אבל לא באופן אחיד.
\item \textbf{מידע חסר:} ג' מקבלת 5, כי כל סוכן חייב להחזיר ציון כיסוי; בגישת הסוכן היחיד הניתוח הוא ברמת הכתבה, בלי מדד לחוסר \cite{zhao2025}; גם בחלופה א' אין מנגנון שיטתי.
\item \textbf{חוזקות וסיכונים:} ג' מקבלת 4, כי יש סוכן לכל ממד, אבל איכות נמדדת באותות עקיפים בלבד.
\item \textbf{אבטחת מידע:} א' מקבלת 5, כי אינה שולחת דבר למודל שפה; ב' 2, כי כל הבקשה, ההקשר והמסקנות עוברים במודל חיצוני אחד; ג' 3, כי היא שולחת למודל רק את שם הספק והרכיב, השער מחושב בקוד בלי מודל, וכל סוכן יכול לעבור בנפרד למודל מקומי.
\item \textbf{ישימות:} ב' מקבלת 5 (פרומפט אחד); א' 4 (בלי פיתוח, אבל מילוי ידני לכל מקרה); ג' 3 (שבעה סוכנים, שער ושקלול).
\item \textbf{גמישות:} ג' מקבלת 5, כי אפשר להחליף סוכן או מודל בלי לגעת בשאר; ב' 3, כי כל שינוי הוא שינוי בפרומפט ארוך אחד; א' 2, כי כל שינוי דורש הדרכה מחדש.
\item \textbf{עלות:} א' מקבלת 2 (שעות אנליסט); ב' 5 (קריאה אחת); ג' 4 (שבע קריאות לפחות, אך לפי \en{AlMahri} ועמיתיו כ-8 סנטים לניתוח \cite{almahri2026}).
\end{itemize}

\textbf{ניתוח רגישות.} בתרחיש שבו אבטחת המידע חשובה הרבה יותר, משקלה עולה ל-0.30 ומשקל ההסבריות יורד ל-0.05; גם אז ג' מובילה, עם 3.95 מול 3.40 לחלופה א' ו-2.80 לחלופה ב'. כשמעלים משקל של קריטריון אחד ומקטינים את השאר באופן יחסי, חלופה א' עוקפת את ג' רק כשמשקל אבטחת המידע עולה על כ-0.43, וחלופה ב' עוקפת את ג' רק כשמשקל הישימות עולה על כ-0.45. גם אם "טיפול במידע חסר" מאבד את כל משקלו, ג' מובילה (4.06). \textbf{ההחלטה: חלופה ג', מערכת מרובת סוכנים עם שערים דטרמיניסטיים.}

\section{ארכיטקטורת המערכת}\label{sec:architecture}

איור~\ref{fig:arch} מציג את הזרימה מבקשת הבדיקה ועד ההמלצה. העיקרון המרכזי הוא \textbf{הפרדה בין מה שמחושב בקוד לבין מה שמודל השפה מעריך}: בדיקת הסנקציות, השקלול וכללי ההחלטה מחושבים בקוד, ולכן ניתנים לבדיקה וחוזרים על עצמם; מודל השפה משמש רק לקריאת טקסט ולהערכה בתוך כל ממד; וההחלטה הסופית נשארת אצל אדם (\en{REQ-NFR-02}).

\begin{figure}[htbp]
\centering
\begin{tikzpicture}[
  box/.style={draw, rounded corners, align=center, font=\footnotesize, minimum height=9mm, inner sep=3pt},
  code/.style={box, fill=gray!15},
  llm/.style={box, fill=blue!8},
  sgA/.style={box, fill=teal!15, text width=1.42cm},
  rkA/.style={box, fill=orange!20, text width=1.42cm},
  human/.style={box, fill=green!12},
  arr/.style={-{Stealth[length=2mm]}, thick},
  node distance=4mm and 5mm]
\node[human, text width=2.8cm] (in) {\foreignlanguage{hebrew}{בקשה: שם הספק, מדינה ורכיב}};
\node[code, text width=2.8cm, below=of in] (er) {\foreignlanguage{hebrew}{זיהוי ישות: שם רשמי, בעלים ורמת מידע \en{A/B/C}}};
\node[code, text width=2.8cm, below=of er] (gate) {\foreignlanguage{hebrew}{שער: סנקציות, בקרת יצוא ומדינה עוינת}};
\node[box, fill=red!12, text width=2.2cm, left=12mm of gate] (rej) {\foreignlanguage{hebrew}{פסילה בשער}};
\node[rkA, below=10mm of gate, xshift=-57mm] (a7) {\foreignlanguage{hebrew}{סוכן סייבר (פסיבי)}};
\node[rkA, right=2mm of a7] (a6) {\foreignlanguage{hebrew}{סוכן אירועים (פתוח)}};
\node[rkA, right=2mm of a6] (a5) {\foreignlanguage{hebrew}{סוכן גיאו\-פוליטי}};
\node[sgA, right=2mm of a5] (a4) {\foreignlanguage{hebrew}{סוכן פיננסי}};
\node[sgA, right=2mm of a4] (a3) {\foreignlanguage{hebrew}{סוכן עמידה ביעדים}};
\node[sgA, right=2mm of a3] (a2) {\foreignlanguage{hebrew}{סוכן התאמה טכנית}};
\node[sgA, right=2mm of a2] (a1) {\foreignlanguage{hebrew}{סוכן איכות ואמינות}};
\node[code, text width=7.4cm, below=10mm of a4] (agg) {\foreignlanguage{hebrew}{שקלול: פרופיל חוזקות $P$, דגלי סיכון $R$, כיסוי $C$, כללי חומרה ושער סייבר מותנה}};
\node[llm, text width=7.4cm, below=of agg] (rep) {\foreignlanguage{hebrew}{כותב הדוח: חוזקות וסיכונים לכל ממד, מקור לכל טענה ורשימת בדיקה ידנית}};
\node[human, text width=7.4cm, below=of rep] (hum) {\foreignlanguage{hebrew}{בודק אנושי: ההחלטה הסופית}};
\draw[arr] (in) -- (er);
\draw[arr] (er) -- (gate);
\draw[arr] (gate) -- (rej);
\foreach \n in {a1,a2,a3,a4,a5,a6,a7}{ \draw[arr] (gate.south) -- (\n.north); }
\foreach \n in {a1,a2,a3,a4,a5,a6,a7}{ \draw[arr] (\n.south) -- (agg.north); }
\draw[arr] (agg) -- (rep);
\draw[arr] (rep) -- (hum);
\end{tikzpicture}

\vspace{1mm}
{\footnotesize \colorbox{teal!15}{סוכן חוזקות} \ \colorbox{orange!20}{סוכן סיכון} \ \colorbox{gray!15}{קוד} \ \colorbox{blue!8}{מודל שפה} \ \colorbox{green!12}{אדם}}
\caption{ארכיטקטורת מצפן ספקים: שער חסימה, שבעה סוכנים במקביל (ארבעה לחוזקות ושלושה לסיכונים), שקלול בקוד והחלטה אנושית}\label{fig:arch}
\end{figure}

\textbf{הפלט של כל סוכן.} נסמן ב-$d$ ממד כלשהו מבין השבעה. הסוכן של ממד $d$ מחזיר ציון $s_d$ בין 0 ל-100 (בשני החלקים ציון גבוה הוא טוב: ספק חזק, או סיכון נמוך), ציון כיסוי $c_d$ בין 0 ל-1 (החלק המשוקלל מהקטגוריות בממד שנבדק בפועל), רשימת טענות עם מקור לכל אחת, ורשימת נושאים שלא נמצא עליהם מידע. טענה בלי מקור אינה נכנסת לציון. סוכן האירועים מוסיף לכל אירוע רמת חומרה (1 עד 4) ורמת ביטחון, ואירוע שאינו מתאים לאף קטגוריה מסומן "אחר / חדש" ומחכה לאישור אדם.

\textbf{השקלול.} $D^{+}$ הם ממדי \textbf{פרופיל החוזקות} (איכות ואמינות, התאמה טכנית, עמידה ביעדים, מצב פיננסי), ו-$D^{-}$ ממדי \textbf{דגלי הסיכון} (מצב גיאופוליטי, אירועים חריגים, אבטחת סייבר). $D^{*}$ היא קבוצת הממדים שבהם $c_d \ge c_{\min}$. המשקלות $w_d$ נקבעים ב-\en{AHP} \textbf{בתוך כל חלק}, וסכומם 1 בכל חלק:
\begin{equation}\label{eq:score}
P=\frac{\sum_{d\in D^{+}\cap D^{*}} w_d\, s_d}{\sum_{d\in D^{+}\cap D^{*}} w_d},\qquad
R=\frac{\sum_{d\in D^{-}\cap D^{*}} w_d\, s_d}{\sum_{d\in D^{-}\cap D^{*}} w_d},\qquad
C=\tfrac{1}{2}\Bigl(\sum_{d\in D^{+}} w_d\, c_d+\sum_{d\in D^{-}} w_d\, c_d\Bigr)
\end{equation}
$R$ הוא ציון דגלי הסיכון, וציון גבוה בו פירושו מעט דגלים. שני חלקים נפרדים מונעים מאיכות גבוהה "לפצות" על סיכון גבוה, ומקטינים את ה-\en{AHP} מ-21 השוואות בזוגות (שבעה ממדים יחד) ל-9 (מטריצה של ארבעה ממדים ומטריצה של שלושה). המכנים מחלקים מחדש, באופן גלוי בדוח, את משקל הממדים שאין עליהם מספיק מידע, ו-$C$ מונע מהיעדר מידע להתפרש כהיעדר סיכון או כאיכות טובה.

\textbf{כללי ההחלטה,} לפי הסדר:
\begin{enumerate}
\item ספק ברשימת סנקציות או בקרת יצוא, או שהוא או בעלי השליטה בו ממדינה עוינת: \textbf{לא מומלץ (פסילה)}, לפני כל חישוב.
\item $C<0.4$, או חלק שאין בו אף ממד עם מספיק מידע: \textbf{מידע לא מספיק}, עם רשימת בדיקה ידנית ושאלון לספק. אם בכל זאת נמצא אירוע ברמה 2 שלא נסגר, ההמלצה היא "\textbf{לא מומלץ}" לפי כלל 4, והאירוע מוצג בראש הדוח.
\item לפי החלק החלש, $m=\min(P,R)$: אם $m\ge 70$ \textbf{מומלץ}; אם $50\le m<70$ \textbf{מומלץ בתנאים}; אם $m<50$ \textbf{לא מומלץ}.
\item \textbf{תקרה:} אירוע ברמת חומרה 2 (כשל אספקה) שלא נסגר מוריד את ההמלצה ל"\textbf{לא מומלץ}" עד לבדיקה ידנית. אירוע סייבר חמור שלא נסגר ולא עצר ייצור מוריד "מומלץ" ל"\textbf{מומלץ בתנאים}" עד לבדיקת אבטחה ידנית. המלצה נמוכה יותר נשארת כמו שהיא.
\end{enumerate}
\textbf{רמת הביטחון} גבוהה אם $C\ge 0.75$ ובינונית אם $0.4\le C<0.75$. הערכים $c_{\min}$, 0.4, 50, 70 ו-0.75 הם ערכי פתיחה: הם נקבעים סופית על \textbf{קבוצת כיול} מתוך מאגר התיקוף, ו-\en{H1} ו-\en{H2} נבדקות על \textbf{קבוצת בדיקה} נפרדת. אם המאגר קטן מכדי לחלק אותו, הספים נשארים ערכי הפתיחה ומדווח ניתוח רגישות שלהם.

\section{הדאטה: מקורות מידע ציבוריים}\label{sec:data}

כל המקורות ציבוריים (\en{REQ-CON-01}), והאיסוף פסיבי בלבד, בלי סריקת מערכות הספק (\en{\mbox{REQ-CON-02}}). טבלה~\ref{tab:data} מפרטת את המקורות לכל ממד. לפני השימוש בכל מקור נבדקים תנאי השימוש, אופן הגישה (ממשק תכנות, \en{API}, או הורדה) והכיסוי הגיאוגרפי (משימה \en{T1.2}, פרק~\ref{ch:workplan}).

{\small
\begin{longtable}{|>{\raggedright\arraybackslash}p{2.2cm}|>{\raggedright\arraybackslash}p{5.8cm}|>{\raggedright\arraybackslash}p{3.8cm}|>{\raggedright\arraybackslash}p{1.6cm}|}
\caption{מקורות מידע ציבוריים מועמדים לכל ממד}\label{tab:data}\\
\hline
\textbf{ממד} & \textbf{מקור מועמד} & \textbf{מה הוא נותן} & \textbf{סוג החברה} \\ \hline
\endfirsthead
\hline
\textbf{ממד} & \textbf{מקור מועמד} & \textbf{מה הוא נותן} & \textbf{סוג החברה} \\ \hline
\endhead
שער חסימה & \en{SDN} של \en{OFAC} \cite{ofacsdn}, \en{Entity List} של \en{BIS} \cite{bis2025}, רשימות האיחוד האירופי, \en{OpenSanctions}, ורשימת המדינות העוינות לפי ההגדרה בפרק~\ref{ch:problem}, מול מקור רשמי שייקבע במשימה \en{T2.1} & האם הישות או בעליה ברשימה, והאם מדינת המושב עוינת או בלי יחסים עם ישראל & כל ספק \\ \hline
גיאופוליטי & מדד \en{GPR} \cite{caldara2022}, מדדי הממשל של הבנק העולמי (\en{WGI}), מרשמי חברות & סיכון מדינה, יחסים עם ישראל ובעלות & ציבורית ופרטית \\ \hline
פיננסי & \en{SEC EDGAR} (\en{10-K}, \en{10-Q}), דוחות בבורסות אחרות, מרשמי חברות שמחייבים דיווח & יחסים פיננסיים ומדד \en{Z} \cite{altman1968}; בפרטית: מדדים עקיפים & ציבורית; פרטית חלקית \\ \hline
עמידה ביעדים & תחזיות בהודעות ובדוחות מול הביצוע & הפער בין ההבטחה לתוצאה & בעיקר ציבורית \\ \hline
אירועים חריגים & חדשות (למשל \en{GDELT}), הודעות רגולטורים, רשומות בתי משפט & אירועים שליליים עם תאריך ומקור & ציבורית ופרטית \\ \hline
סייבר & דיווחי \en{8-K} סעיף 1.05 \cite{sec2023}, הודעות \en{CERT} ורגולטורים, מאגרי \en{ISO 27001} & אירועים שדווחו והסמכות & ציבורית ופרטית \\ \hline
איכות ואמינות & \en{OASIS} של \en{IAQG} (\en{AS9100}), \en{USAspending.gov}, פרסי ספקים, הוראות תקינות (\en{Airworthiness Directives}) & הסמכות, לקוחות, הכרה פומבית, ריקולים & ציבורית ופרטית \\ \hline
התאמה טכנית & מפרטים באתר החברה, פטנטים, הסמכות תהליך (\en{Nadcap}) & התאמת חומרים, תהליכים ויכולות לרכיב & ציבורית ופרטית \\* \hline
\end{longtable}
}

\textbf{מאגר התיקוף.} המתודולוגיה תיבדק לאחור על מקרים ציבוריים לפי סולם החומרה (פרק~\ref{ch:problem}): ברמה 2, \en{Incora} (פשיטת רגל, 2023) \cite{incora2023} ו-\en{ASCO} (כופרה, 2019) \cite{certeu2019}; ברמה 1, חברות שנוספו לרשימות סנקציות ובקרת יצוא בגלל העברת רכיבי תעופה לאיראן \cite{treasury2023,bis2025}; וברמה 3, \en{FACC}, יצרנית אוסטרית של חלקי מטוסים שאיבדה 42 מיליון אירו ב-2016 בהונאת "מנכ"ל מזויף" \cite{reuters2016facc}, כדי לבדוק שאות סיכון אינו מסווג ככשל. מולם תיבחר קבוצה של 5 עד 10 ספקים מוכחים, שמהם מצופה "מומלץ". בניית המאגר המלא היא משימה \en{T1.1}.

לכל מקרה נכתבים מראש התוצאה הצפויה ותאריכי החיתוך. \textbf{כשל שאפשר לצפות}, כמו \en{Incora}, שקדמו לו ירידה בהכנסות ונטל חוב, משמש לבדיקת החיזוי ולמדידת \textbf{זמן ההתרעה}: מספר החודשים בין תאריך החיתוך המוקדם ביותר שבו ההמלצה היא "לא מומלץ" לבין מועד הכשל (שאלת משנה 4). זה אותו רף שלפיו נקבע הסיווג הנכון בפרק~\ref{ch:problem}; התאריך שבו ההמלצה ירדה לראשונה מ"מומלץ" מדווח בנפרד, כאות מוקדם בלבד. \textbf{כשל שאי אפשר לצפות}, כמו מתקפת הכופרה על \en{ASCO}, משמש רק לבדיקת סיווג החומרה. טבלה~\ref{tab:cases} מפרטת לכל סוג מקרה מתי הוא מורץ, מה התוצאה הצפויה, והאם הוא נכנס לחישוב הסיווג הנכון של \en{H1}.

{\small
\begin{longtable}{|>{\raggedright\arraybackslash}p{3.3cm}|>{\raggedright\arraybackslash}p{3.6cm}|>{\raggedright\arraybackslash}p{3.9cm}|>{\raggedright\arraybackslash}p{3.3cm}|}
\caption{התוצאה הצפויה לכל סוג מקרה במאגר התיקוף}\label{tab:cases}\\
\hline
\textbf{המקרה} & \textbf{תאריך החיתוך} & \textbf{התוצאה הצפויה} & \textbf{חישוב \en{H1}} \\ \hline
\endfirsthead
\hline
\textbf{המקרה} & \textbf{תאריך החיתוך} & \textbf{התוצאה הצפויה} & \textbf{חישוב \en{H1}} \\ \hline
\endhead
\en{Incora}: כשל שאפשר לצפות & 12 חודשים לפני פשיטת הרגל, ועוד שני תאריכים מוקדמים יותר & "לא מומלץ" & כן; נמדד גם זמן ההתרעה \\ \hline
\en{ASCO}: כשל שאי אפשר לצפות & אחרי פרסום המתקפה & אירוע ברמת חומרה 2, ולפי כלל 4 "לא מומלץ" & לא; בודק את סיווג החומרה (\en{\mbox{REQ-FUN-07}}) \\ \hline
ישויות ברשימות סנקציות ובקרת יצוא & אחרי הכניסה לרשימה & "פסילה" בשער & לא; בודק את שער החסימה (\en{\mbox{REQ-FUN-01}}) \\ \hline
\en{FACC}: אות סיכון & אחרי פרסום ההונאה & אירוע ברמת חומרה 3; ההמלצה אינה יורדת בגללו ל"לא מומלץ" & לא; בודק שאות סיכון אינו מסווג ככשל \\ \hline
ספקים מוכחים & מועד ההרצה & "מומלץ" & כן \\* \hline
\end{longtable}
}

\textbf{מניעת "ידיעת העתיד".} מודל שפה כבר יודע ש-\en{Incora} פשטה רגל. לכן כל הרצה נעשית \textbf{לפי נקודת זמן}: הסוכנים מקבלים רק מסמכים מלפני תאריך החיתוך, ההנחיה אוסרת ידע מחוץ למסמכים, וכל מקור שצוטט נבדק מול תאריך החיתוך. המקרים שלעיל קדמו לתאריך הידע של המודלים, ולכן זו ההגנה היחידה עליהם; מקרים מ-2025 ואילך יתווספו אם יימצאו (סיכון \en{\mbox{VAL-02}}, פרק~\ref{ch:risks}).

\section{החדשנות והערך המוסף}\label{sec:innovation}

החדשנות אינה במערכת ניטור חדשה, אלא בשלוש נקודות שלא נמצאו בספרות (פרק~\ref{ch:literature}) ובתיעוד המוצרים (פרק~\ref{ch:market}):
\begin{enumerate}
\item \textbf{הערכה בשני חלקים עם ציון כיסוי:} ההמלצה עונה על "כמה הספק טוב", "כמה הוא מסוכן" ו"כמה ידוע עליו", כך שספק פרטי לא נראה בטוח רק כי לא נמצאו עליו חדשות רעות.
\item \textbf{חומרה לפי רציפות האספקה, עם קטגוריה פתוחה:} אירוע מדורג לפי השפעתו על האספקה, ואירוע מסוג חדש אינו נכפה על טקסונומיה קיימת.
\item \textbf{שערים בקוד, הערכה במודל, החלטה אצל אדם:} כל המלצה ניתנת לבדיקה, ואפשר להעביר את המערכת למודל מקומי ברשת הארגון בלי לשנות את הלוגיקה (\en{REQ-CON-03}).
\end{enumerate}

% ---------- פרק 6: תוכנית עבודה ----------
% לפי פרק 8 במדריך. הספירה מתחילה ב-11.11.2026; אבן דרך 1 = הגשת הדוח המכין ב-25.11.2026 (עמית, 16.11).
% RACI לפי עמית (16.11): מנחה תמיד A; תע"א C במשימות 1,3,6,8 ו-I במשימות 7,9. 22.11: עמודות בשמות (עמית סוסן, טל אזואלוס); R/R ב-T0, T1.3, T5.1 נשאר לבקשת עמית.
\chapter{תוכנית עבודה}\label{ch:workplan}

הפרק מתאר איך המתודולוגיה תמומש, ועונה על שלוש שאלות: מה עושים, מי עושה ומתי. לוח הזמנים מתחיל ב-11.11.2026 ומתפרס על 16 שבועות, עם שלוש אבני דרך:
\begin{enumerate}
\item \textbf{הגשת הדוח המכין} ב-25.11.2026, בסוף שבוע 2.
\item \textbf{אב-טיפוס עובד} בסוף שבוע 10, לפי יעד ה-\en{SMART} בפרק~\ref{ch:problem}.
\item \textbf{השלמת התיקוף} בסוף שבוע 13.
\end{enumerate}

\section{פירוק העבודה (\en{WBS})}\label{sec:wbs}

הפרויקט מחולק לחמש חבילות עבודה (\en{Work Packages}). לכל חבילה יש משימות ותוצר מוגדר (איור~\ref{fig:wbs} וטבלה~\ref{tab:wbs}).

\begin{figure}[!htbp]
\centering
\begin{tikzpicture}[
  wp/.style={draw, rounded corners, fill=blue!6, align=center, font=\footnotesize, text width=2.4cm, minimum height=12mm},
  top/.style={draw, rounded corners, fill=gray!15, align=center, font=\small\bfseries, text width=4.5cm, minimum height=9mm}]
\node[top] (p) at (0,0) {\foreignlanguage{hebrew}{אב-טיפוס מצפן ספקים}};
\foreach \i/\t [count=\k] in {
  5/{תיעוד והגנה},
  4/{תיקוף},
  3/{שקלול ודוח},
  2/{שער וסוכנים},
  1/{נתונים והכנה}}{
  \node[wp] (w\i) at ({(\k-3)*2.8},-2) {\en{WP\i}\\\foreignlanguage{hebrew}{\t}};
  \draw[thick] (p.south) -- ++(0,-0.4) -| (w\i.north);
}
\end{tikzpicture}
\caption{מבנה פירוק העבודה: חמש חבילות תחת פרויקט אחד}\label{fig:wbs}
\end{figure}

{\small
\begin{longtable}{|>{\raggedright\arraybackslash}p{1.3cm}|>{\raggedright\arraybackslash}p{2.6cm}|>{\raggedright\arraybackslash}p{5.6cm}|>{\raggedright\arraybackslash}p{4.3cm}|}
\caption{חבילות העבודה, המשימות והתוצרים}\label{tab:wbs}\\
\hline
\textbf{מזהה} & \textbf{החבילה} & \textbf{משימות} & \textbf{תוצר} \\ \hline
\endfirsthead
\hline
\textbf{מזהה} & \textbf{החבילה} & \textbf{משימות} & \textbf{תוצר} \\ \hline
\endhead
\en{WP1} & נתונים והכנה & \en{T1.1} בניית מאגר התיקוף: ספקים שכשלו וספקים מוכחים, עם תוצאה צפויה לכל מקרה (טבלה~\ref{tab:cases}). \en{T1.2} חיבור למקורות המידע ובדיקת תנאי השימוש שלהם. \en{T1.3} מילוי \en{AHP} וניתוח רגישות & מאגר תיקוף מתועד, מחברים למקורות, משקלות סופיים \\ \hline
\en{WP2} & שער וסוכנים & \en{T2.1} זיהוי הישות ושער הסנקציות. \en{T2.2} שבעה סוכני ממד: ארבעה לפרופיל החוזקות ושלושה לדגלי הסיכון, עם פלט במבנה קבוע & קוד מתועד ובדיקות יחידה \\ \hline
\en{WP3} & שקלול ודוח & \en{T3.1} מנוע השקלול, כללי ההחלטה וכותב הדוח & אב-טיפוס שמפיק דוח מלא \\ \hline
\en{WP4} & תיקוף & \en{T4.1} בדיקה לאחור לפי נקודת זמן על ספקים שכשלו ועל ספקים מוכחים, עם שני מודלי שפה לפחות, וכיול הספים על קבוצת הכיול. \en{T4.2} שאלון משוב מהתעשייה האווירית (\en{H3}) & דוח תיקוף: \en{H1} עד \en{H3} \\ \hline
\en{WP5} & תיעוד והגנה & \en{T5.1} ספר הפרויקט, מצגת והגנה & הגשה סופית \\ \hline
\end{longtable}
}

\section{מטריצת אחריות (\en{RACI})}\label{sec:raci}

טבלה~\ref{tab:raci} מגדירה לכל משימה את התפקיד של כל משתתף: \textbf{\en{R}} (\en{Responsible}) מבצעת את המשימה ואחראית לתוצר; \textbf{\en{A}} (\en{Accountable}) מאשר את התוצר; \textbf{\en{C}} (\en{Consulted}) מתייעצים איתו במהלך הביצוע; \textbf{\en{I}} (\en{Informed}) מקבל עדכון על התוצאה.

{\small
\begin{longtable}{|c|>{\raggedright\arraybackslash}p{5.6cm}|c|c|c|c|}
\caption{מטריצת אחריות}\label{tab:raci}\\
\hline
\textbf{\#} & \textbf{משימה} & \textbf{עמית סוסן} & \textbf{טל אזואלוס} & \textbf{מנחה אקדמי} & \textbf{תע"א} \\ \hline
\endfirsthead
\hline
\textbf{\#} & \textbf{משימה} & \textbf{עמית סוסן} & \textbf{טל אזואלוס} & \textbf{מנחה אקדמי} & \textbf{תע"א} \\ \hline
\endhead
0 & \en{T0} כתיבת הדוח המכין & \en{R} & \en{R} & \en{A} & --- \\ \hline
1 & \en{T1.1} מאגר תיקוף: ספקים שכשלו וספקים מוכחים & \en{R} & \en{C} & \en{A} & \en{C} \\ \hline
2 & \en{T1.2} מחברים למקורות מידע & \en{C} & \en{R} & \en{A} & --- \\ \hline
3 & \en{T1.3} \en{AHP} וניתוח רגישות & \en{R} & \en{R} & \en{A} & \en{C} \\ \hline
4 & \en{T2.1} זיהוי ישות ושער סנקציות & \en{C} & \en{R} & \en{A} & --- \\ \hline
5 & \en{T2.2} סוכני הממדים: חוזקות וסיכונים & \en{R} & \en{C} & \en{A} & --- \\ \hline
6 & \en{T3.1} שקלול ודוח & \en{C} & \en{R} & \en{A} & \en{C} \\ \hline
7 & \en{T4.1} בדיקה לאחור וכיול & \en{R} & \en{C} & \en{A} & \en{I} \\ \hline
8 & \en{T4.2} שאלון משוב & \en{C} & \en{R} & \en{A} & \en{C} \\ \hline
9 & \en{T5.1} ספר, מצגת והגנה & \en{R} & \en{R} & \en{A} & \en{I} \\ \hline
\end{longtable}
}

\section{לוח זמנים ותרשים גאנט}\label{sec:gantt}

טבלה~\ref{tab:tasks} מפרטת את המשימות, את התלויות ביניהן ואת אבני הדרך. איור~\ref{fig:gantt} מציג אותן בתרשים גאנט. שבוע 1 מתחיל ב-11.11.2026.

\textbf{הנתיב הקריטי} הוא שרשרת המשימות שבה כל עיכוב מעכב את כל הפרויקט:
\begin{center}
\en{T0 $\rightarrow$ T1.2 $\rightarrow$ T2.1 $\rightarrow$ T2.2 $\rightarrow$ T3.1 $\rightarrow$ T4.1 $\rightarrow$ T4.2 $\rightarrow$ \foreignlanguage{hebrew}{רזרבה} $\rightarrow$ T5.1}
\end{center}
אורך הנתיב הוא \en{$2+3+1+3+1+2+1+1+2=16$} שבועות.

ל-\en{T1.1} ול-\en{T1.3} יש מרווח: הן צריכות להסתיים רק לפני \en{T4.1} ו-\en{T3.1}. כך, עיכוב קצר בבניית מאגר התיקוף לא יעכב את הפרויקט.

\begin{figure}[!htbp]
\centering
\begin{tikzpicture}[x=0.66cm, y=0.46cm]
% ציר שבועות
\foreach \w in {1,...,16}{
  \node[font=\scriptsize] at (\w-0.5,0.7) {\en{\w}};
  \draw[gray!25] (\w,0.3) -- (\w,-11.3);
}
\draw[gray!25] (0,0.3) -- (0,-11.3);
% משימות: שורה, שם, התחלה, סוף, קריטית?
\foreach \row/\name/\s/\e/\c in {
  1/T0/1/2/1,
  2/T1.1/3/4/0,
  3/T1.2/3/5/1,
  4/T1.3/4/5/0,
  5/T2.1/6/6/1,
  6/T2.2/7/9/1,
  7/T3.1/10/10/1,
  8/T4.1/11/12/1,
  9/T4.2/13/13/1,
  10/{\foreignlanguage{hebrew}{רזרבה}}/14/14/1,
  11/T5.1/15/16/1}{
  \node[font=\scriptsize, anchor=east] at (-0.2,-\row+0.5) {\name};
  \ifnum\c=1
    \fill[red!65] (\s-1,-\row+0.15) rectangle (\e,-\row+0.85);
  \else
    \fill[blue!45] (\s-1,-\row+0.15) rectangle (\e,-\row+0.85);
  \fi
}
% אבני דרך
\foreach \x/\y/\n in {2/-0.5/1, 10/-6.5/2, 13/-8.5/3}{
  \fill[green!50!black] (\x,\y) +(0,0.3) -- +(0.3,0) -- +(0,-0.3) -- +(-0.3,0) -- cycle;
  \node[font=\scriptsize, anchor=west] at (\x+0.3,\y) {\foreignlanguage{hebrew}{אבן דרך \n}};
}
% מקרא
\fill[red!65] (0,-12.3) rectangle (0.6,-11.9); \node[font=\scriptsize, anchor=west] at (0.7,-12.1) {\foreignlanguage{hebrew}{נתיב קריטי}};
\fill[blue!45] (5,-12.3) rectangle (5.6,-11.9); \node[font=\scriptsize, anchor=west] at (5.7,-12.1) {\foreignlanguage{hebrew}{משימה עם מרווח}};
\end{tikzpicture}
\caption{תרשים גאנט ל-16 שבועות, החל מ-11.11.2026. הנתיב הקריטי באדום, ושלוש אבני הדרך בירוק.}\label{fig:gantt}
\end{figure}

{\small
\begin{longtable}{|l|>{\raggedright\arraybackslash}p{5cm}|c|>{\raggedright\arraybackslash}p{2.2cm}|>{\raggedright\arraybackslash}p{3.4cm}|}
\caption{משימות, תלויות ואבני דרך}\label{tab:tasks}\\
\hline
\textbf{קוד} & \textbf{המשימה} & \textbf{שבועות} & \textbf{קודמות} & \textbf{אבן דרך / תוצר} \\ \hline
\endfirsthead
\hline
\textbf{קוד} & \textbf{המשימה} & \textbf{שבועות} & \textbf{קודמות} & \textbf{אבן דרך / תוצר} \\ \hline
\endhead
\en{T0} & כתיבת הדוח המכין & \en{1--2} & --- & \textbf{1: הגשת הדוח המכין (25.11.2026)} \\ \hline
\en{T1.1} & מאגר תיקוף: ספקים שכשלו וספקים מוכחים & \en{3--4} & \en{T0} & מאגר תיקוף \\ \hline
\en{T1.2} & מחברים למקורות מידע ובדיקת תנאי שימוש & \en{3--5} & \en{T0} & גישה לכל מקור \\ \hline
\en{T1.3} & \en{AHP} וניתוח רגישות & \en{4--5} & \en{T0} & משקלות סופיים \\ \hline
\en{T2.1} & זיהוי ישות ושער סנקציות & 6 & \en{T1.2} & שער עובד \\ \hline
\en{T2.2} & שבעה סוכני ממד & \en{7--9} & \en{T2.1} & פלט לכל ממד \\ \hline
\en{T3.1} & שקלול, כללי החלטה ודוח & 10 & \en{T2.2}, \en{T1.3} & \textbf{2: אב-טיפוס עובד} \\ \hline
\en{T4.1} & בדיקה לאחור בשני מודלים וכיול ספים & \en{11--12} & \en{T3.1}, \en{T1.1} & תוצאות \en{H1}, \en{H2} \\ \hline
\en{T4.2} & שאלון משוב (\en{H3}) & 13 & \en{T4.1} & \textbf{3: התיקוף הושלם} \\ \hline
--- & רזרבה & 14 & --- & --- \\ \hline
\en{T5.1} & ספר פרויקט, מצגת והגנה & \en{15--16} & \en{T4.2} & הגשה והגנה \\ \hline
\end{longtable}
}

\textbf{מעקב.} כל פגישה עם המנחה מתועדת בפרוטוקול קצר: מה הושלם, מה התעכב ומה הצעד הבא. הקוד והדוח מנוהלים במאגר \en{GitHub}, וכך אפשר לראות בכל רגע את מצב העבודה מול התוכנית.

% ---------- פרק 7: ניתוח סיכונים ותוכנית הפחתה ----------
% 22.11.2026: גוף שלישי, קיצור, איור בעברית. ערכי L ו-I הם הערכה איכותנית. ORG-04 (איש קשר בתע"א) הוסר לבקשת המרצה, 16.11.
\chapter{ניתוח סיכונים ותוכנית הפחתה}\label{ch:risks}

לכל סיכון מוערכות \textbf{הסתברות} $L$ ו\textbf{השפעה} $I$ בסולם 1 עד 5 (בהערכה איכותנית), וציון הסיכון הוא $E = L \times I$. הסימונים שונים מ-$P$ ו-$R$ של פרק~\ref{ch:design}. סיכון עם $E\ge 12$ הוא \textbf{קריטי}, ולכל סיכון נקבעות \textbf{פעולת מנע}, שננקטת מראש, ו\textbf{תוכנית מגירה}, שמופעלת אם הוא מתממש.

\section{זיהוי הסיכונים וכימותם}\label{sec:risk-id}

טבלה~\ref{tab:risks} מפרטת את הסיכונים, ואיור~\ref{fig:riskmatrix} ממקם אותם במטריצה. מבין ארבעת הסיכונים הקריטיים, שניים נוגעים לתוקף התיקוף (\en{DAT-01}, \en{VAL-02}) ושניים לאמינות הפלט של המערכת (\en{ALG-03}, \en{ENT-07}).

{\small
\begin{longtable}{|l|>{\raggedright\arraybackslash}p{7.2cm}|>{\raggedright\arraybackslash}p{1.9cm}|c|c|c|}
\caption{זיהוי הסיכונים וכימותם}\label{tab:risks}\\
\hline
\textbf{מזהה} & \textbf{תרחיש הכשל} & \textbf{תחום} & $L$ & $I$ & $E$ \\ \hline
\endfirsthead
\hline
\textbf{מזהה} & \textbf{תרחיש הכשל} & \textbf{תחום} & $L$ & $I$ & $E$ \\ \hline
\endhead
\en{DAT-01} & אין מספיק מידע ציבורי על מקרי התיקוף ועל אותות איכות, במיוחד בחברות פרטיות, ולכן אי אפשר לבדוק את \en{H1} ו-\en{H2} & נתונים & 4 & 4 & \textbf{16} \\ \hline
\en{VAL-02} & "ידיעת העתיד": המודל כבר יודע מה קרה לספק, והבדיקה לאחור נראית טובה מהמציאות & תיקוף & 4 & 4 & \textbf{16} \\ \hline
\en{ALG-03} & הזיה: הסוכן ממציא עובדה או מקור, והם נכנסים להמלצה & אלגוריתם & 3 & 5 & \textbf{15} \\ \hline
\en{SEC-04} & מידע פנימי נשלח בטעות לשירות ענן, או ששם הספק והרכיב חושפים כוונת רכש & אבטחת מידע & 2 & 5 & 10 \\ \hline
\en{LEG-05} & תנאי השימוש של מקור מידע אוסרים איסוף אוטומטי & משפטי & 3 & 3 & 9 \\ \hline
\en{CST-06} & עלות הקריאות למודל או מגבלות קצב חורגות מהתקציב & משאבים & 2 & 2 & 4 \\ \hline
\en{ENT-07} & זיהוי ישות שגוי: שם דומה, חברה-בת או שינוי שם גורמים לבדיקת הישות הלא נכונה, והשער מחמיץ ספק ברשימה & נתונים & 3 & 4 & \textbf{12} \\ \hline
\end{longtable}
}

\begin{figure}[!htbp]
\centering
\begin{tikzpicture}[x=0.95cm,y=0.6cm]
\foreach \p in {1,...,5}{
 \foreach \i in {1,...,5}{
   \pgfmathtruncatemacro{\r}{\p*\i}
   \ifnum\r>11 \def\col{red!35}\else\ifnum\r>5 \def\col{yellow!40}\else\def\col{green!20}\fi\fi
   \fill[\col] (\i-1,\p-1) rectangle (\i,\p);
   \draw[gray!50] (\i-1,\p-1) rectangle (\i,\p);
   \node[font=\tiny, gray] at (\i-0.15,\p-0.82) {\r};
 }}
\foreach \i in {1,...,5}{ \node[font=\scriptsize] at (\i-0.5,-0.3) {\i}; \node[font=\scriptsize] at (-0.3,\i-0.5) {\i}; }
\node[font=\small] at (2.5,-0.9) {\foreignlanguage{hebrew}{השפעה} $I$};
\node[font=\small, rotate=90] at (-0.9,2.5) {\foreignlanguage{hebrew}{הסתברות} $L$};
\node[font=\scriptsize, align=center] at (3.5,3.5) {DAT-01\\VAL-02};
\node[font=\scriptsize] at (4.5,2.5) {ALG-03};
\node[font=\scriptsize] at (3.5,2.5) {ENT-07};
\node[font=\scriptsize] at (4.5,1.5) {SEC-04};
\node[font=\scriptsize] at (2.5,2.5) {LEG-05};
\node[font=\scriptsize] at (1.5,1.5) {CST-06};
\end{tikzpicture}
\caption{מטריצת סיכונים: הסתברות מול השפעה. באדום הסיכונים הקריטיים ($E\ge 12$)}\label{fig:riskmatrix}
\end{figure}

\section{פעולות מנע ותוכניות מגירה}\label{sec:mitigation}

טבלה~\ref{tab:mitigation} מגדירה לכל סיכון פעולת מנע, תוכנית מגירה וציון שיורי אחרי פעולת המנע.

{\small
\begin{longtable}{|l|>{\raggedright\arraybackslash}p{6.4cm}|>{\raggedright\arraybackslash}p{4.4cm}|c|}
\caption{פעולות מנע, תוכניות מגירה וציון שיורי}\label{tab:mitigation}\\
\hline
\textbf{מזהה} & \textbf{פעולת מנע} & \textbf{תוכנית מגירה} & \textbf{שיורי} \\ \hline
\endfirsthead
\hline
\textbf{מזהה} & \textbf{פעולת מנע} & \textbf{תוכנית מגירה} & \textbf{שיורי} \\ \hline
\endhead
\en{DAT-01} & בחירת מקרים עם ציר זמן מתועד בכמה מקורות (\en{T1.1}); "מידע לא מספיק" הוא תוצאה נכונה שגם היא נבדקת & הצגת התוצאות כמקרי בוחן תיאוריים & 8 \\ \hline
\en{VAL-02} & הרצה לפי נקודת זמן, הנחיה שאוסרת ידע מחוץ למסמכים, ובדיקת תאריך כל מקור שצוטט & העדפת מקרים מאחרי תאריך הידע, ודיווח על המגבלה & 8 \\ \hline
\en{ALG-03} & טענה בלי מקור אינה נכנסת לציון; בדיקה אוטומטית שהקישור קיים והציטוט מופיע בו; השערים בקוד. הבדיקה אינה מבטיחה שהמסקנה נובעת מהמקור, ולכן הסיכון השיורי אינו מינימלי & ביקורת ידנית על מדגם טענות (\en{REQ-NFR-01}), והחלפת מודל & 10 \\ \hline
\en{SEC-04} & רק מידע ציבורי נכנס למערכת; שמות רכיבים כלליים בתיקוף; ארכיטקטורה שמאפשרת מודל מקומי (\en{REQ-CON-03}) & עצירת ההרצה, מחיקת המידע ודיווח למנחה & 5 \\ \hline
\en{LEG-05} & בדיקת תנאי השימוש לפני חיבור לכל מקור (\en{T1.2}); עדיפות לממשקי תכנות רשמיים & מקור חלופי מאותו סוג (טבלה~\ref{tab:data}) & 4 \\ \hline
\en{CST-06} & מדידת העלות לדוח באב-הטיפוס, ומודל קטן במשימות פשוטות, שבמאמר הנלווה כמעט לא נפל בביצועים \cite{zhao2025} & הגבלת מספר הספקים בהרצה & 2 \\ \hline
\en{ENT-07} & זיהוי הישות בקוד (\en{T2.1}) מול מרשם עם מזהה ייחודי, כולל חברות-בת ושמות קודמים, ואישור המשתמש לפני ההרצה & סימון "זיהוי לא ודאי" והעברה לבדיקה ידנית & 8 \\ \hline
\end{longtable}
}

אחרי פעולות המנע אף סיכון אינו נשאר קריטי. הגבוה ביותר הוא \en{ALG-03} (10), כי בדיקת הקישור אינה מבטיחה שהמסקנה נובעת מהמקור; אחריו \en{DAT-01}, \en{VAL-02} ו-\en{ENT-07} (8). \en{DAT-01} ו-\en{VAL-02} נוגעים לתוקף התיקוף, ולכן תוצאותיו יוצגו כמקרי בוחן ולא כהוכחה סטטיסטית.

% ---------- פרק 8: סיכום ----------
% לפי פרק 10 במדריך: ממצאים, מוכנות (Design Readiness) ושלוש פעולות מיידיות אחרי אישור הדוח. עמוד אחד.
\chapter{סיכום}\label{ch:summary}

בארגונים גדולים קליטת ספק חדש נמשכת שבועות \cite{hicx2021}, הסיכון הפיננסי גבוה יותר דווקא בחברות פרטיות, שעליהן יש פחות מידע \cite{rapidratings2026}, וגם ארגון ביטחוני גדול אינו יכול להסתמך רק על מה שהספק מוסר \cite{gao2025}. מאמר הדגל \cite{sadeghi2026} מראה שאפשר לשלב מודל שפה בבחירת ספקים רב-קריטריונית, והמאמר הנלווה \cite{zhao2025} מראה שמודלי שפה מזהים סיכוני ספקים מתוך חדשות, אבל אף אחד מהם לא מפיק המלצה על ספק בודד מתוך מידע ציבורי ולא מודד כמה מהמידע חסר (פרק~\ref{ch:literature}). גם בתיעוד הפומבי של המוצרים המסחריים לא נמצא מוצר שמעריך את איכות הספק, מציג ציון כיסוי מידע ומקשר כל טענה למקור (פרק~\ref{ch:market}). מבין שלוש חלופות נבחרה, במטריצת החלטה משוקללת ובניתוח רגישות, מערכת מרובת סוכנים שמעריכה כל ספק בשני חלקים, פרופיל חוזקות ודגלי סיכון, ומחשבת את שער הפסילה בקוד (פרק~\ref{ch:design}).

\section*{מוכנות לשלב המימוש}
הארכיטקטורה, הנוסחאות, מקורות המידע ותוכנית העבודה מוגדרים, ולכל הסיכונים הקריטיים נקבעו פעולות מנע ותוכניות מגירה (פרקים~\ref{ch:workplan} ו-\ref{ch:risks}). המשקלות הסופיים ומאגר התיקוף נקבעים במשימות הראשונות של תוכנית העבודה (\en{T1.1} ו-\en{T1.3}), ולכן אינם מעכבים את תחילת המימוש.

\section*{פעולות מיידיות}
אחרי אישור הדוח יבוצעו שלוש פעולות (איור~\ref{fig:next}): פגישת התנעה עם המנחה האקדמי והמנחה התעשייתי לאישור תוכנית העבודה, פתיחת מאגר \en{GitHub} לקוד, לנתונים ולפרוטוקולי הפגישות, ותחילת המשימות הראשונות בתוכנית העבודה, מילוי ה-\en{AHP} ובניית מאגר התיקוף.

\begin{figure}[!h]
\centering
\begin{tikzpicture}[
  st/.style={draw, rounded corners, fill=blue!6, align=center, font=\footnotesize, text width=3.1cm, minimum height=10mm},
  arr/.style={-{Stealth[length=2mm]}, thick}]
\node[st] (c) at (0,0) {\foreignlanguage{hebrew}{מילוי ה-\en{AHP} ובניית מאגר התיקוף}};
\node[st] (b) at (3.8,0) {\foreignlanguage{hebrew}{פתיחת מאגר \en{GitHub} לקוד ולנתונים}};
\node[st] (a) at (7.6,0) {\foreignlanguage{hebrew}{פגישת התנעה עם שני המנחים}};
\node[st, fill=green!12] (ok) at (11.4,0) {\foreignlanguage{hebrew}{הדוח אושר}};
\draw[arr] (ok) -- (a); \draw[arr] (a) -- (b); \draw[arr] (b) -- (c);
\end{tikzpicture}
\caption{שלוש הפעולות המיידיות אחרי אישור הדוח}\label{fig:next}
\end{figure}


% ---------- מקורות ----------
\clearpage
% השורה בתוכן העניינים נכתבת כאן, בהקשר עברי, כדי שלא תתהפך
\addcontentsline{toc}{chapter}{מקורות}
% ---------- רשימת מקורות (סגנון IEEE, ממוספרת לפי סדר ההופעה בטקסט) ----------
% כל מקור כאן מופיע גם ב"מקורות - רשימה מאומתת.md" עם סטטוס האימות שלו.
% מוסיפים מקור חדש לפי הסדר שבו הוא מצוטט לראשונה בדוח.
\begin{otherlanguage}{english}
\renewcommand{\bibname}{\foreignlanguage{hebrew}{מקורות}}
\raggedright
\small\setstretch{1.15}
\begin{thebibliography}{99}

\bibitem{rapidratings2026}
S. Samora, ``Production demand raises financial stress on aerospace, defense suppliers: Report,'' \emph{Manufacturing Dive}, May 19, 2026. [Online]. Available: \url{https://manufacturingdive.com/news/defense-aerospace-suppliers-rapid-ratings-financial-risks-james-gellert/819966}

\bibitem{gao2025}
U.S. Government Accountability Office, ``Defense Industrial Base: Actions Needed to Address Risks Posed by Dependence on Foreign Suppliers,'' GAO-25-107283, Jul. 24, 2025. [Online]. Available: \url{https://files.gao.gov/reports/GAO-25-107283/index.html}

\bibitem{enisa2021}
European Union Agency for Cybersecurity (ENISA), ``Threat Landscape for Supply Chain Attacks,'' Jul. 2021. [Online]. Available: \url{https://www.enisa.europa.eu/publications/threat-landscape-for-supply-chain-attacks}

\bibitem{certeu2019}
CERT-EU, ``Threat Memo 190617-1'' (TLP:WHITE), Jun. 17, 2019. [Online]. Available: \url{https://cert.europa.eu/static/MEMO/2019/TLP-WHITE-CERT-EU-MEMO-190617-1.pdf}

\bibitem{cmmc2024}
U.S. Department of Defense, ``Cybersecurity Maturity Model Certification (CMMC) Program,'' final rule, 32 CFR Part 170, \emph{Federal Register}, vol. 89, no. 199, p. 83092, doc. 2024-22905, Oct. 15, 2024 (effective Dec. 16, 2024). [Online]. Available: \url{https://www.govinfo.gov/content/pkg/FR-2024-10-15/html/2024-22905.htm}

\bibitem{zhao2025}
M. Zhao, O. Hussain, Y. Zhang, M. Saberi, and A. Leshob, ``Benchmarking large language models for supply chain risk identification: An extended evaluation within the LARD-SC framework,'' \emph{Service Oriented Computing and Applications}, 2025, doi: 10.1007/s11761-025-00474-7.

\bibitem{almahri2026}
S. AlMahri, L. Xu, and A. Brintrup, ``Automating supply chain disruption monitoring via an agentic AI approach,'' arXiv:2601.09680, Jan. 2026 (preprint).

\bibitem{mevaker2024}
State Comptroller of Israel, ``Prevention of fraud and embezzlement at Israel Aerospace Industries Ltd.,'' in \emph{Report of the State Comptroller of Israel: The Defense System}, May 2024, abstract, p. 59. [Online]. Available: \url{https://library.mevaker.gov.il/sites/DigitalLibrary/Documents/2024/2024.05/EN/2024-05-105-IAI-Taktzir-EN.pdf}

\bibitem{iai2026}
Israel Aerospace Industries, ``IAI annual financial results for 2025,'' press release, Mar. 11, 2026. [Online]. Available: \url{https://www.iai.co.il/press/iai-annual-financial-results-for-2025/}

\bibitem{hicx2021}
HICX, ``The Supplier Experience Survey 2021,'' Jun. 2021. [Online]. Available: \url{https://hicx.com/wp-content/uploads/2021/06/0043_ARTICLE-The-Supplier-Experience-Survey-2021-v7.pdf}

\bibitem{incora2023}
B. Unglesbee, ``Incora's bankruptcy sends a warning signal to aerospace supply chains,'' \emph{Supply Chain Dive}, Jun. 9, 2023. [Online]. Available: \url{https://www.supplychaindive.com/news/incoras-bankruptcy-aerospace-supply-chains/652433}

% ---- נוסף בפרק 2 (14.11.2026) ----
\bibitem{sec2023}
U.S. Securities and Exchange Commission, ``SEC adopts rules on cybersecurity risk management, strategy, governance, and incident disclosure by public companies,'' press release 2023-139, Jul. 26, 2023. [Online]. Available: \url{https://www.sec.gov/newsroom/press-releases/2023-139}

\bibitem{altman1968}
E. I. Altman, ``Financial ratios, discriminant analysis and the prediction of corporate bankruptcy,'' \emph{The Journal of Finance}, vol. 23, no. 4, pp. 589--609, 1968, doi: 10.1111/j.1540-6261.1968.tb00843.x.

% ---- נוספו בפרק 3 (סקירת ספרות), לפי סדר ההופעה ----
\bibitem{ho2010}
W. Ho, X. Xu, and P. K. Dey, ``Multi-criteria decision making approaches for supplier evaluation and selection: A literature review,'' \emph{European Journal of Operational Research}, vol. 202, no. 1, pp. 16--24, 2010, doi: 10.1016/j.ejor.2009.05.009.

\bibitem{saaty1990}
T. L. Saaty, ``How to make a decision: The analytic hierarchy process,'' \emph{European Journal of Operational Research}, vol. 48, no. 1, pp. 9--26, 1990, doi: 10.1016/0377-2217(90)90057-I.

\bibitem{boukrouh2024}
I. Boukrouh, F. Tayalati, and A. Azmani, ``A comprehensive framework for supplier selection: Using subjective, objective, and hybrid multi-criteria decision-making techniques with sensitivity analysis,'' \emph{IEEE Access}, vol. 12, pp. 145550--145569, 2024, doi: 10.1109/ACCESS.2024.3462348.

\bibitem{dey2015}
P. K. Dey, A. Bhattacharya, and W. Ho, ``Strategic supplier performance evaluation: A case-based action research of a UK manufacturing organisation,'' \emph{International Journal of Production Economics}, vol. 166, pp. 192--214, 2015, doi: 10.1016/j.ijpe.2014.09.021.

\bibitem{ayyadi2024}
I. Ayyadi and B. Oulhadj, ``ISO 9001 certification and organizational performance among manufacturing companies: A systematic literature review and research agenda,'' \emph{Business Strategy \& Development}, vol. 7, no. 4, Art. no. e70044, 2024, doi: 10.1002/bsd2.70044.

\bibitem{ho2015}
W. Ho, T. Zheng, H. Yildiz, and S. Talluri, ``Supply chain risk management: A literature review,'' \emph{International Journal of Production Research}, vol. 53, no. 16, pp. 5031--5069, 2015, doi: 10.1080/00207543.2015.1030467.

\bibitem{caldara2022}
D. Caldara and M. Iacoviello, ``Measuring geopolitical risk,'' \emph{American Economic Review}, vol. 112, no. 4, pp. 1194--1225, 2022, doi: 10.1257/aer.20191823.

\bibitem{bai2025}
X. Bai, K. Z. Zhou, B. Zhang, and E. Xie, ``The dual effects of geopolitical risk on MNCs' first-tier supply base: A political economy perspective,'' \emph{Journal of International Business Studies}, vol. 56, no. 9, pp. 1125--1149, 2025, doi: 10.1057/s41267-025-00805-y.

\bibitem{charpin2025}
R. Charpin and M. Cousineau, ``Friendshoring: How geopolitical tensions affect foreign sourcing, supply base complexity, and sub-tier supplier sharing,'' \emph{International Journal of Operations \& Production Management}, vol. 45, no. 5, pp. 1006--1031, 2025, doi: 10.1108/IJOPM-01-2024-0067.

\bibitem{zhaoj2024}
J. Zhao, J. Ouenniche, and J. De Smedt, ``Survey, classification and critical analysis of the literature on corporate bankruptcy and financial distress prediction,'' \emph{Machine Learning with Applications}, vol. 15, Art. no. 100527, 2024, doi: 10.1016/j.mlwa.2024.100527.

\bibitem{chernakov2025}
P. Chernakov, S. Jafarnejad, and R. Frank, ``An agentic LLM framework for adverse media screening in AML compliance,'' arXiv:2602.23373, 2025 (preprint).

\bibitem{zhang2026}
L. Zhang, X. Sun, M. Xu, Y. Yang, T. Zhang, and Z. Wen, ``Large language model-driven analysis of supply risk dynamics in global electric vehicle supply chain,'' \emph{Environmental Impact Assessment Review}, vol. 119, Art. no. 108378, 2026, doi: 10.1016/j.eiar.2026.108378.

\bibitem{nist2022}
J. Boyens, A. Smith, N. Bartol, K. Winkler, A. Holbrook, and M. Fallon, ``Cybersecurity supply chain risk management practices for systems and organizations,'' NIST Special Publication 800-161 Rev.~1 (update 1), National Institute of Standards and Technology, Nov. 2024, doi: 10.6028/NIST.SP.800-161r1-upd1.

\bibitem{gani2023}
A. B. D. Gani, Y. Fernando, S. Lan, M. K. Lim, and M.-L. Tseng, ``Interplay between cyber supply chain risk management practices and cyber security performance,'' \emph{Industrial Management \& Data Systems}, vol. 123, no. 3, pp. 843--861, 2023, doi: 10.1108/IMDS-05-2022-0313.

\bibitem{slapnicar2025}
S. Slapni{\latinfallback č}ar, T. Vidmar, and E. Tsen, ``Process theory of supplier cyber risk assessment,'' \emph{Australasian Journal of Information Systems}, vol. 29, 2025, doi: 10.3127/ajis.v29.5323.

\bibitem{jackson2024}
I. Jackson, D. Ivanov, A. Dolgui, and J. Namdar, ``Generative artificial intelligence in supply chain and operations management: A capability-based framework for analysis and implementation,'' \emph{International Journal of Production Research}, vol. 62, no. 17, pp. 6120--6145, 2024, doi: 10.1080/00207543.2024.2309309.

\bibitem{huang2025}
L. Huang \emph{et al.}, ``A survey on hallucination in large language models: Principles, taxonomy, challenges, and open questions,'' \emph{ACM Transactions on Information Systems}, vol. 43, no. 2, pp. 1--55, 2025, doi: 10.1145/3703155.

\bibitem{sadeghi2026}
N. Sadeghi, M. Dehghani, and N. Kabadayi, ``The AI-driven decision-making (AIDM) framework: Integrating AHP and ChatGPT for supplier selection,'' \emph{Annals of Operations Research}, vol. 359, pp. 2593--2615, 2026, doi: 10.1007/s10479-026-07136-7.

% ---- נוספו בפרק 4 (חקר שוק) ----
\bibitem{wiseguy2026}
WiseGuy Reports, ``Third party \& supplier risk management software market,'' report WGR642606, 2026 (commercial market report). [Online]. Available: \url{https://www.wiseguyreports.com/reports/third-party-supplier-risk-management-software-market}

\bibitem{everstream2026}
Everstream Analytics, ``Everstream Analytics named a Leader in the 2026 Gartner Magic Quadrant for Supplier Risk Management Solutions,'' press release, May 7, 2026 (cites Gartner, C. Curtin, M. Shreffler, and C. Van Dyke, May 5, 2026). [Online]. Available: \url{https://www.everstream.ai/media/everstream-analytics-named-a-leader-in-the-2026-gartner-magic-quadrant-for-supplier-risk-management-solutions/}

\bibitem{sipri2025}
L. Scarazzato, N. Tian, D. Lopes da Silva, X. Liang, Z. A. Karim, and J. Guiberteau Ricard, ``The SIPRI Top 100 arms-producing and military services companies, 2024,'' SIPRI Fact Sheet, Dec. 2025. [Online]. Available: \url{https://www.sipri.org/publications/2025/sipri-fact-sheets/sipri-top-100-arms-producing-and-military-services-companies-2024}

\bibitem{rrfaq}
RapidRatings, ``Financial health FAQs,'' vendor website, accessed Nov. 13, 2026. [Online]. Available: \url{https://rapidratings.com/financial-health-intelligence/financial-health-faqs-enterprise}

\bibitem{interosgov}
Interos, ``Government supply chains,'' vendor website, accessed Nov. 13, 2026. [Online]. Available: \url{https://www.interos.ai/industries/federal-government}

\bibitem{interosnavy2023}
``US Navy selects Interos for supply chain risk management,'' \emph{Security Journal Americas}, Jun. 6, 2023. [Online]. Available: \url{https://securityjournalamericas.com/us-navy-interos-supply-chain-risk/}

\bibitem{bitsight2026}
Bitsight, ``How does Bitsight work? A look at security ratings \& how they're used,'' vendor blog, accessed Nov. 13, 2026. [Online]. Available: \url{https://www.bitsight.com/blog/how-does-bitsight-work}

\bibitem{prewave2026}
Prewave, ``Gartner Magic Quadrant,'' vendor website, accessed Nov. 13, 2026. [Online]. Available: \url{https://www.prewave.com/gartner-magic-quadrant}

% ---- נוספו בפרק 5 (תכן ודאטה) ----
\bibitem{ofacsdn}
U.S. Office of Foreign Assets Control, ``Specially Designated Nationals and Blocked Persons List (SDN),'' sanctions list site, accessed Nov. 16, 2026. [Online]. Available: \url{https://sanctionslist.ofac.treas.gov/Home/SdnList}

\bibitem{bis2025}
U.S. Bureau of Industry and Security, ``Additions to the Entity List,'' final rule, 15 CFR Part 744, \emph{Federal Register}, vol. 90, no. 194, p. 48193, doc. 2025-19508, Oct. 9, 2025. [Online]. Available: \url{https://www.govinfo.gov/content/pkg/FR-2025-10-09/pdf/2025-19508.pdf}

\bibitem{treasury2023}
U.S. Department of the Treasury, ``Treasury sanctions multinational network supporting Iran's UAV and military aircraft production,'' press release JY1745, Sep. 19, 2023. [Online]. Available: \url{https://home.treasury.gov/news/press-releases/jy1745}

\bibitem{reuters2016facc}
Reuters, ``Austria's FACC, hit by cyber fraud, fires CEO,'' May 25, 2016 (republished by Yahoo Finance). [Online]. Available: \url{https://ca.finance.yahoo.com/news/austrias-facc-hit-cyber-fraud-091143816.html}


\end{thebibliography}
\end{otherlanguage}


\end{document}
